\clearpage
\appendix

\section{Supplementary Material}
\label{sec:appendix}

Unless noted otherwise, all analyses use OpenVLA-7B (\texttt{openvla/openvla-7b} on HuggingFace) evaluated on LIBERO-Goal (10 tasks $\times$ 50 demonstrations, $N = 63{,}228$ observations after excluding terminal frames).
Model conditions, representation extraction, metric definitions, and probe architectures follow \S\ref{sec:method}.
Cross-suite experiments (\S\ref{sec:cross_benchmark_appendix}) use LIBERO-Spatial with the same protocol (10 tasks $\times$ 50 demonstrations).
All experiments ran on 2$\times$ NVIDIA RTX PRO 6000 Blackwell Server Edition GPUs (96\,GB each); random seed is 42 unless stated otherwise.

\subsection{Architecture Summary}
Architecture details are provided in Table~\ref{tab:arch_overview} (main text).

\subsection{OpenVLA-Specific Analysis Methods}
\label{sec:openvla_methods_appendix}

Three additional analyses are applied to OpenVLA to characterize instruction encoding and representational geometry.

\paragraph{Cross-instruction probes.}
Each visual frame is paired with all 10 LIBERO-Goal task instructions, producing 10 representation vectors per frame per layer.
A 10-way linear classifier identifies which instruction was present, applied separately at image\_mean and readout, with leave-one-out cross-validation (LOOCV) over 10 folds.

\paragraph{Representational similarity analysis (RSA).}
RSA~\cite{kriegeskorte2008} quantifies geometric alignment between representations and the action space.
For each task, we compute pairwise representational dissimilarity matrices (RDMs) using cosine distances for representations and Euclidean distances for 7-D action vectors, from within-task observation pairs across 50 demonstrations.
Alignment is measured via Spearman $\rho$ between upper-triangular RDM entries; we report mean $\rho$ across 10 tasks with 95\% confidence intervals from 10{,}000 bootstrap resamples at the task level~\cite{nili2014toolbox}.
We prefer RSA over CKA~\cite{kornblith2019}: RSA captures relative distance structure via scale-invariant rank correlation, suited to comparing high- and low-dimensional spaces (4{,}096-D vs.\ 7-D), while CKA measures global subspace overlap that can remain high when geometric structure differs.
Both metrics are reported in \S\ref{sec:cka_appendix}; their divergence (\S\ref{sec:routing_synthesis}) confirms that neither alone fully characterizes representational structure.

\paragraph{External baselines.}
DINOv2 ViT-L/14~\cite{oquab2024_dinov2} (\texttt{cls\_token}/\texttt{patch\_mean}) and CLIP ViT-L/14~\cite{radford2021_clip} (\texttt{patch\_mean}) provide self-supervised and vision-language pretrained reference points for probing and RSA.

\subsection{OpenVLA Single-Architecture Deep Dive}
\label{sec:openvla_deepdive_appendix}

This section provides the full single-architecture analysis summarized in \S\ref{sec:openvla_routing}.

\paragraph{Cross-instruction probes.}
Cross-instruction probes validate the analysis framework on OpenVLA~\cite{kim2024_openvla}: as predicted by the causal mask, image\_mean instruction classification stays at chance ($10\%$) across all 32 layers in all conditions (full profile in \S\ref{sec:dense_routing}).
Trained readout achieves $100\%$ 10-way classification from L1 onward and maintains this through L31 (LOOCV-confirmed; \S\ref{sec:loocv}).
Reference readout also reaches $100\%$ at L1 but decays to $82\%$ by L31; the readout token is thus the primary multimodal integration site, and instruction encoding remains stable across depth in the trained model.
All patterns reproduce with a linear probe (\S\ref{sec:linear_probe}).

\paragraph{L1 attention gradient.}
At L1, three conditions form a gradient of instruction attention: randomly initialized networks allocate 0/32 instruction-dominant heads ($11.7\%$ text attention, $1.5\times$ text-to-image ratio), Llama-2 base~\cite{touvron2023_llama2} allocates 9/32 ($36.5\%$), and VLA-trained allocates 24/32 ($57.7\%$, $38.2\times$ text-to-image ratio; \S\ref{sec:attention_decomp}).
The 32-layer attention profile reveals qualitatively distinct organizations: Llama-2 base concentrates instruction attention at mid-depth (L9--L16: 15--22 heads), while VLA fine-tuning produces an entry-exit bookend pattern with L0--L1 at 21--24 heads, L31 at 18, and intermediate layers at 0--11 (\S\ref{sec:threshold_sensitivity}).
Head counts are computed at a $50\%$ threshold; sensitivity analysis ($30\%$--$70\%$) confirms ordering robustness.
A fourth condition, OpenVLA fine-tuned on LIBERO, preserves L1 entry concentration (22/32 heads) but loses L31 exit re-engagement entirely (0 vs.\ 18 heads for the OXE-pretrained model), consistent with exit-layer attention serving multi-task flexibility rather than task-specific routing.
The redistribution beyond Llama-2 base priors is consistent with VLA fine-tuning being associated with instruction routing beyond what is inherited from language-model attention.

\paragraph{Counterfactual perturbation.}
Four perturbation types across 10 LIBERO-Goal tasks produce large action divergence: wrong task ($d = 1.182$), empty instruction ($1.170$), shuffled tokens ($1.096$), paraphrased ($0.996$).
All four conditions yield $d \geq 1.0$, with ``wrong'' slightly exceeding ``empty.''
The effect replicates across LIBERO-Spatial (10 tasks): empty $1.173$, shuffled $0.979$, wrong $0.885$, paraphrased $0.807$ (Table~\ref{tab:counterfactual_spatial}).
Across both suites (20 tasks), the gradient spans $d = 0.81$--$1.18$.
Spatial instructions differ only in prepositions, yet wrong-task perturbation still produces $d = 0.885$; instruction-dominant heads route fine-grained spatial information, not merely task identity.

\paragraph{Dose-response ablation.}
Dose-response ablation of L1 instruction-dominant heads produces monotonic trajectory divergence: $d$ rises from $0.482$ to $0.639$ as 4--32 heads are progressively ablated, with steep increase through 12 heads and plateau beyond 20 (\S\ref{sec:action_trajectory}).
Ablating all 24 instruction-dominant heads alters $72.3\%$ of action tokens (token-level $d = 0.65$, trajectory-level $d = 2.943$); control ablation of 8 random heads yields $d = 0.41$ ($29.3\%$ tokens affected), concentrated on the gripper dimension.
The routing effect is task-dependent: instruction-discriminative tasks show $2.5\times$ larger $|\Delta R^2|$ under L1 ablation than visually-sufficient tasks ($p = 0.033$; \S\ref{sec:task_complexity}).

\paragraph{$R^2$ localization.}
$R^2$ action probing localizes training effects: trained readout peaks at $R^2 = 0.678$ (L13), reference readout peaks at $0.609$ (L32), and the training effect concentrates at the readout ($+0.069$) rather than at image\_mean ($+0.015$; \S\ref{sec:action_r2_threecond}).
These effects survive Bonferroni correction across all 63 tests (\S\ref{sec:bonferroni_correction}).

\subsection{Counterfactual Sensitivity Details}
\label{sec:cf_details_appendix}

\begin{table}[h]
\caption{\textbf{Counterfactual sensitivity summary.} Cohen's $d$ range and most sensitive perturbation for action-producing models (OpenVLA, TraceVLA). NORA uses cosine distance on representations (different scale; \S\ref{sec:cf_repr}).}
\label{tab:cf_summary}
\small
\centering
\setlength{\tabcolsep}{5pt}
\begin{tabular}{@{}llcl@{}}
\toprule
\textbf{Architecture} & \textbf{$d$ range} & \textbf{Most sensitive} & \textbf{Metric} \\
\midrule
OpenVLA & 0.45--1.21 & empty (2/4) & Cohen's $d$ \\
TraceVLA & 0.97--1.52 & empty (4/4) & Cohen's $d$ \\
\midrule
NORA    & 0.28--sat. & --- & cosine dist.$^{a}$ \\
\bottomrule
\multicolumn{4}{l}{\footnotesize $^{a}$Representation-space only; not comparable to Cohen's $d$.}\\
\multicolumn{4}{l}{\footnotesize Empty ranks first in 6/8 cells across both architectures.}\\
\end{tabular}
\end{table}

Table~\ref{tab:counterfactual} reports per-condition counterfactual rankings for all eight cells where Cohen's $d$ is available.

\begin{table}[h]
\centering
\caption{Counterfactual condition ranking for trained models (Cohen's $d$; higher $=$ greater action divergence under instruction perturbation). Conditions: E $=$ empty instruction, W $=$ wrong task, S $=$ shuffled tokens, P $=$ paraphrased. ``Empty'' ranks first in 6 of 8 cells; OpenVLA Goal (``wrong'' first) and Long (``shuffled'' first) are the exceptions. NORA cells use cosine distance and are reported in \S\ref{sec:cf_sensitivity}.}
\label{tab:counterfactual}
\small
\begin{tabular}{llcccc}
\toprule
\textbf{Arch.} & \textbf{Suite} & \textbf{1st} & \textbf{2nd} & \textbf{3rd} & \textbf{4th} \\
\midrule
OpenVLA & Goal    & W\,1.18 & E\,1.17 & S\,1.10 & P\,1.00 \\
OpenVLA & Spatial & E\,1.17 & S\,0.98 & W\,0.89 & P\,0.81 \\
OpenVLA & Object  & E\,1.21 & W\,1.19 & S\,1.12 & P\,0.90 \\
OpenVLA & Long    & S\,0.57 & W\,0.55 & E\,0.51 & P\,0.45 \\
\midrule
TraceVLA & Goal    & E\,1.52 & W\,1.40 & S\,1.34 & P\,1.21 \\
TraceVLA & Spatial & E\,1.33 & S\,1.30 & P\,1.18 & W\,1.17 \\
TraceVLA & Object  & E\,1.21 & S\,1.16 & W\,1.06 & P\,0.97 \\
TraceVLA & Long    & E\,1.47 & W\,1.42 & S\,1.39 & P\,1.16 \\
\bottomrule
\end{tabular}
\end{table}

\paragraph{Metric constraints on trained--reference comparison.}
Trained models produce action outputs, enabling Cohen's $d$ on action MSE (OpenVLA, TraceVLA) or cosine distance on action tokens (NORA).
Reference VLMs (Phi-3-Vision, Qwen2.5-VL) do not generate actions; their sensitivity is measured via cosine distance on hidden representations, a metric with different scale, units, and distributional properties.
Ratios between Cohen's $d$ and cosine distance are therefore methodologically invalid, and we do not report cross-metric amplification factors.
Qualitatively, reference VLMs show weak sensitivity: the TraceVLA reference (Phi-3-Vision) produces peak cosine distances of $0.15$--$0.19$; the NORA reference (Qwen2.5-VL) yields $0.081$ on Goal.

\paragraph{NORA details.}
NORA Goal produces saturated exit-layer cosine distances of $1.0$ for all conditions, with differentiation visible only at earlier layers (L24: wrong $= 0.123$, empty $= 0.098$).
NORA Spatial produces a peak cosine distance of $0.275$ for the ``wrong'' condition.

\subsection{CKA Analysis Method}
\label{sec:cka_appendix}

We employ Centered Kernel Alignment (CKA)~\cite{kornblith2019} alongside RSA as a complementary perspective.
While RSA captures local geometric alignment via rank-order preservation, CKA measures global representational similarity through kernel-based comparisons.

We compute both \emph{linear CKA}, which uses a linear kernel and is sensitive to second-order statistics, and \emph{RDM kernel CKA}, which operates on the same distance matrices as RSA but aggregates them differently.
Using both variants tests whether conclusions about representation structure are robust to the choice of similarity metric.

Linear CKA between high-dimensional representations (4{,}096-D for LLM layers) and the 7-D action space is subject to a known dimension mismatch: the linear kernel in the low-dimensional space may not capture the same structure as in the high-dimensional space.
RDM kernel CKA avoids this issue by operating on distance matrices of equal size regardless of the original feature dimensionality.

\subsection{Metric Sensitivity: RSA versus CKA}
Linear CKA~\cite{kornblith2019} reverses the RSA ordering: reference LLM ($\text{CKA} = 0.089$) $>$ trained LLM ($0.083$) $>$ DinoSigLIP+Projector ($0.081$) $>$ DINOv2 ($0.066$), partially inverting the self-supervised advantage (Figure~\ref{fig:metric_comparison}).
RDM kernel CKA partially resolves this discrepancy: trained ($0.260$) $\approx$ Projector ($0.257$) $>$ reference ($0.246$) $>$ DINOv2 ($0.241$).
This divergence reflects the known sensitivity of representation similarity metrics to different geometric properties.

At the per-layer level, linear CKA between image\_mean and the action space rises monotonically from L0 ($0.081$) to L32 ($0.098$), with no inverted-U pattern.
RSA captures rank-order geometric alignment; CKA captures global linear structure.
The dissociation indicates that the geometric refinement in early layers operates on the relative arrangement of representations (affecting RSA) without altering the overall covariance structure (captured by CKA).

\begin{figure}[t]
\centering
\includegraphics[width=\columnwidth]{figures/fig6.pdf}
\caption{Three similarity metrics across representation sources. \textbf{(a)}~RSA ranks DINOv2 cls ($0.692$) $>$ Fused ($0.634$) $>$ Projector ($0.615$). \textbf{(b)}~Linear CKA reverses the ordering: Reference ($0.089$) $>$ Trained ($0.083$) $>$ Projector ($0.081$) $>$ DINOv2 ($0.066$). \textbf{(c)}~RDM kernel CKA: Trained ($0.260$) $\approx$ Projector ($0.257$) $>$ Reference ($0.246$) $>$ DINOv2 ($0.241$). The metric divergence underscores that RSA and CKA capture different aspects of representational structure.}
\label{fig:metric_comparison}
\end{figure}

\subsection{Branch-Split Detailed Analysis}
To disentangle the contributions of the two encoder streams within DinoSigLIP, we extract representations from the DINOv2 branch (1{,}024-D) and SigLIP branch (1{,}152-D) separately before fusion.
The two branches yield near-identical RSA: DINOv2 branch $\rho = 0.619$, SigLIP branch $\rho = 0.621$.
Fusion produces a representation ($\rho = 0.634$) that exceeds either branch alone; concatenation captures complementary structure.
The MLP projector then compresses this signal: DinoSigLIP+Projector achieves $\rho = 0.615$, a drop of $\Delta = 0.019$ from the fused 2{,}176-D input.

\subsection{RSA Across Representation Sources}
Figure~\ref{fig:rsa_bar} shows RSA with the action space across all representation sources; the ordering matches the RSA results reported in \S\ref{sec:routing_synthesis}.

\begin{figure}[h]
\centering
\includegraphics[width=\columnwidth]{figures/fig2.pdf}
\caption{RSA with the action space ($\rho$) across representation sources. Self-supervised DINOv2 \texttt{cls\_token} ($\rho{=}0.692$) achieves the highest geometric alignment. CLIP ($0.557$) yields the weakest alignment among pretrained encoders. Error bars: 95\% bootstrap CIs.}
\label{fig:rsa_bar}
\end{figure}

\subsection{Per-Layer RSA Profiles}
% Figure promoted to main text (Figure~\ref{fig:rsa_perlayer} in \S\ref{sec:cross_arch_analyses}).
The per-layer RSA profile (Figure~\ref{fig:rsa_perlayer} in main text) extends Figure~\ref{fig:r2_collapse} by adding readout and reference profiles.

\subsection{Per-Layer Action Decoding Across Training Stages}
\label{sec:action_r2_threecond}

Figure~\ref{fig:action_r2_threecond} compares per-layer action probe $R^2$ across three training stages (randomly initialized, Llama-2 base, VLA-trained) using 50 demonstrations per task, averaged over 3 seeds (mean $\pm$ std).
All three conditions share L0 representations ($R^2 = 0.645 \pm 0.008$); subsequent differences thus arise from Transformer processing rather than input variation.
The trained model peaks at L9 ($R^2 = 0.643 \pm 0.014$), consistent with progressive specialization for action-relevant feature integration.
Llama-2 base shows a steady decline (L0: $0.645 \to$ L31: $0.557$, std $0.003$--$0.028$), consistent with stable but unspecialized processing.
The randomly initialized model exhibits high cross-seed variance and deep-layer collapse, with negative $R^2$ and bimodal behavior across seeds at late layers; this instability indicates that randomly initialized Transformer layers produce destructive interference on action-relevant structure.
Random weights at intermediate layers produce anti-correlated representations ($R^2 < 0$, L10--L11), consistent with randomly initialized networks lacking structured computation.
The three-condition gradient (trained $>$ Llama-2 base $>$ randomly initialized) holds from L1 through L15, with a reversal at L32 where trained $R^2$ drops below the pretrained baseline.
The qualitative pattern is consistent across all 3 seeds.
Three conditions show substantial agreement in layer-wise $R^2$ ordering (Kendall's $W = 0.72$, $\chi^2 = 68.99$, $p < 0.001$).

% Figure promoted to main text (Figure~\ref{fig:action_r2_threecond} in \S\ref{sec:openvla_routing}).

\subsection{Random-Subset Ablation Details}
To evaluate whether the ablation effect in \S\ref{sec:openvla_routing} is specific to instruction-dominant heads, we compare dominant-head ablation against 30 random subsets of 24/32 L1 heads (each subset drawn without replacement).
Dominant-head ablation increases action $R^2$ by $+0.030$ (within-experiment baseline $0.524$), while random subsets yield a mean increase of $+0.071 \pm 0.013$ (Figure~\ref{fig:random_subset_ablation}).
The dominant-head effect falls below all 30 random subsets ($z = -3.11$, percentile $= 0\%$); heads with higher instruction-attention fraction carry more organized attention signals, consistent with graded rather than binary specialization.

\begin{figure}[h]
\centering
\includegraphics[width=0.75\columnwidth]{figures/fig_specificity_bootstrap_30seeds.pdf}
\caption{$R^2$ increase from ablating 24/32 L1 heads. Gray bars: 30 random subsets; red line: dominant-24 heads ($+0.030$). Random subsets yield larger gains (mean $+0.071 \pm 0.013$; $z = -3.11$, percentile $= 0\%$), consistent with graded attention specialization across L1 heads.}
\label{fig:random_subset_ablation}
\end{figure}

\subsection{Action Trajectory Analysis Details}
\label{sec:action_trajectory}

Table~\ref{tab:action_shift_perdim} reports per-dimension mean absolute shift for instruction-dominant (24 heads) and control (8 heads) L1 ablations.
Both conditions produce negligible shifts on translational dimensions ($dx$, $dy$, $dz$ $< 0.01$) and small rotational shifts ($rx$, $ry$, $rz$ $< 0.04$).
The gripper dimension absorbs the majority of the effect: instruction-head ablation shifts gripper predictions by $0.291$ (control: $0.141$), consistent with L1 attention influencing discrete open/close decisions.

Per-task Cohen's $d$ values for instruction-head ablation are $0.43$, $0.69$, $0.70$, $0.70$, and $0.73$ (mean $0.65$), with all tasks significant ($p < 10^{-10}$).
Token-level KL divergence between original and ablated action distributions is $5.08$ for instruction heads and $1.69$ for control heads.

\paragraph{Trajectory-level aggregation.}
Aggregating per-demo mean actions to the trajectory level reveals that token-level $d$ values are conservative: within-trajectory averaging reduces noise, yielding ${\sim}4\times$ larger effect sizes (mean token-to-trajectory ratio $0.245$).
Instruction-head ablation ($N = 50$: 5 tasks $\times$ 10 demos) yields trajectory $d = 2.943$ (95\% CI $[2.469, 3.707]$); bottom-24 non-instruction heads ($N = 30$) yield $d = 2.347$ (CI $[1.854, 3.221]$).
All 14 conditions produce trajectory $d > 2.0$ with CI lower bounds $> 1.7$.
Counterfactual trajectory-level $d$ (10 LIBERO-Goal tasks, $N = 100$): wrong $3.798$ (CI $[3.194, 4.938]$), empty $3.438$ ($[2.921, 4.448]$), shuffled $2.595$ ($[2.028, 3.791]$), paraphrased $2.131$ ($[1.740, 2.819]$).
Token-level $d$ underestimates trajectory-level effects by ${\times}3.1$--$4.5$, consistent with within-trajectory averaging reducing noise.
Fine-grained ordering is not reliably distinguished at trajectory level ($N = 100$); token-level comparisons ($N {\sim} 5{,}000$) provide higher resolution for ordering claims.

\paragraph{Ground-truth action error.}
Ablation redirects model outputs ($d = 0.65$) but does not systematically increase error relative to expert demonstrations: original MSE $= 0.498$; top-24 instruction-head ablation MSE $= 0.497$ ($\Delta = -0.001$, Wilcoxon $p = 0.52$, n.s.); bottom-24 non-instruction heads MSE $= 0.506$ ($\Delta = +0.008$).
Changes concentrate on the gripper dimension; translation and rotation dimensions are negligibly affected.
Per-task breakdown shows all $|\Delta\text{MSE}| < 3\%$ of baseline with all $p > 0.27$; the near-zero aggregate $\Delta$ reflects uniformly small per-task effects rather than cancellation across tasks.
This is consistent with the information-computation dissociation: instruction-dominant heads control instruction-specific direction selection rather than action precision.

Table~\ref{tab:perhead_bootstrap} reports per-head bootstrap CIs ($10{,}000$ resamples).
Individual control heads produce ${\sim}1.9\times$ stronger effects ($d = 0.051$) than instruction heads ($d = 0.027$, $p = 0.0018$), consistent across all 5 tasks.
The 24 instruction heads nonetheless produce a larger aggregate perturbation ($d = 0.65$ vs.\ $0.41$), consistent with collective coordination rather than individual strength.

\begin{table}[h]
\centering
\caption{Per-head bootstrap analysis of L1 ablation effects. Instruction heads are individually weaker but collectively stronger, consistent with coordinated attention redistribution.}
\label{tab:perhead_bootstrap}
\small
\begin{tabular}{lccc}
\toprule
\textbf{Group} & \textbf{$N$ heads} & \textbf{Collective $d$} & \textbf{Per-head $d$ [95\% CI]} \\
\midrule
Instruction-dominant & 24 & 0.65 & 0.027 [0.026, 0.029] \\
Control & 8 & 0.41 & 0.051 [0.048, 0.055] \\
\midrule
Difference & --- & --- & 0.024, $p = 0.0018$ \\
\bottomrule
\end{tabular}
\end{table}

Instruction-head ablation produces more uniform perturbation across action dimensions (CV${}=1.30$ vs.\ control CV${}=2.74$, permutation $p = 0.016$), consistent with structured, language-conditioned computation rather than diffuse noise.

Table~\ref{tab:dose_response} reports the dose-response data for instruction-head ablation at L1 (Figure~\ref{fig:synthesis_panels} in main text).
Trajectory divergence increases monotonically from $d = 0.482$ (4 heads) to $d = 0.639$ (32 heads), with steep rise through 12 heads, plateau at 16--20, and slow increase at 20--32.

\begin{table}[h]
\centering
\caption{Dose-response data for L1 instruction-head ablation. Cohen's $d$ vs.\ original action sequence; 95\% bootstrap CIs ($10{,}000$ resamples).}
\label{tab:dose_response}
\small
\begin{tabular}{rcc}
\toprule
\textbf{Heads ablated} & \textbf{Cohen's $d$} & \textbf{95\% CI} \\
\midrule
4 & 0.482 & [0.449, 0.511] \\
8 & 0.508 & [0.479, 0.541] \\
12 & 0.560 & [0.529, 0.593] \\
16 & 0.592 & [0.559, 0.625] \\
20 & 0.593 & [0.559, 0.627] \\
24 & 0.611 & [0.578, 0.645] \\
28 & 0.619 & [0.587, 0.654] \\
32 & 0.639 & [0.603, 0.671] \\
\bottomrule
\end{tabular}
\end{table}

\begin{table}[h]
\centering
\caption{Counterfactual instruction analysis across all 10 LIBERO-Goal tasks. Jaccard: token overlap with original instruction; Cohen's $d$: mean across 10 tasks. ``Wrong'' slightly exceeds ``empty'' ($d = 1.182$ vs.\ $1.170$); all four conditions produce large effects ($d \geq 1.0$). L1 ablation included for comparison.}
\label{tab:counterfactual_goal}
\small
\begin{tabular}{lcc}
\toprule
\textbf{Condition} & \textbf{Jaccard} & \textbf{Cohen's $d$} \\
\midrule
Empty (no instruction) & 0.000 & 1.170 \\
Wrong (different task) & 0.174 & 1.182 \\
Shuffled (random order) & 1.000 & 1.096 \\
Paraphrased & 0.533 & 0.996 \\
\midrule
L1 ablation (24 heads) & --- & 0.650 \\
\bottomrule
\end{tabular}
\end{table}

\begin{table}[h]
\centering
\caption{Mean absolute action shift per dimension under L1 head ablation. Instruction: 24 dominant heads; Control: 8 remaining heads. Gripper shift dominates both conditions.}
\label{tab:action_shift_perdim}
\small
\begin{tabular}{lccccccc}
\toprule
\textbf{Condition} & \textbf{$dx$} & \textbf{$dy$} & \textbf{$dz$} & \textbf{$rx$} & \textbf{$ry$} & \textbf{$rz$} & \textbf{Gripper} \\
\midrule
Instruction (24) & 0.002 & 0.005 & 0.008 & 0.021 & 0.031 & 0.034 & 0.291 \\
Control (8) & 0.000 & 0.001 & 0.002 & 0.008 & 0.014 & 0.018 & 0.141 \\
\bottomrule
\end{tabular}
\end{table}

\subsection{Pretrained LLM Representation Quality Under Visual Input}
A potential concern with the three-condition comparison (\S\ref{sec:arch_data}) is that the pretrained Llama-2 backbone never encountered projected visual tokens during pretraining, so its representations of such input could be degenerate.
We evaluate three diagnostics to assess representation quality.
Activation magnitude grows from $16.5$ (L0) to $90.2$ (L31) for Llama-2 base, compared to $16.5 \to 51.7$ for the VLA-trained model; the pretrained LLM produces larger but non-exploding activations with no evidence of collapse.
Per-layer cosine similarity between L0 and deeper layers decreases progressively from $1.0$ to $\approx 0.28$ at L25, consistent with gradual differentiation rather than catastrophic out-of-distribution behavior.
Top-10 principal component variance explained remains stable across all layers ($0.59$--$0.74$), with no low-rank collapse.
These diagnostics indicate that the pretrained LLM produces structured, full-rank representations of visual input despite never encountering projected visual tokens during pretraining (Figure~\ref{fig:ood_diagnostics}). These results support the validity of the three-condition comparison.

\begin{figure}[h]
\centering
\includegraphics[width=\columnwidth]{figures/fig_ood_diagnostics.pdf}
\caption{Pretrained LLM representation diagnostics under visual input. Left: activation magnitude across layers (Llama-2 base vs.\ VLA-trained). Center: cosine similarity between L0 and deeper layers. Right: top-10 PC variance explained. All three metrics indicate structured, non-degenerate representations.}
\label{fig:ood_diagnostics}
\end{figure}

\subsection{Probe Architecture Sensitivity}
\label{sec:linear_probe}

To rule out that our findings depend on MLP probe capacity ($N/\text{params} \approx 0.06$, overparameterized), we replicate the per-layer analysis with a linear probe ($4{,}096 \to 7$, $N/\text{params} \approx 2.2$).
The inverted-U pattern reproduces: trained image\_mean peaks at L14 ($R^2 = 0.565$), while the reference declines monotonically from L0 ($R^2 = 0.517$).
The trained--reference gap is larger under the linear probe (L14: $\Delta = +0.324$) than the MLP probe (L8: $\Delta = +0.015$); trained representations are thus more probe-decodable for action information.
Reference deep layers collapse entirely under the linear probe (L25: $R^2 = -0.298$); MLP nonlinearity partially rescues these degraded representations.
For readout tokens, the linear probe also corroborates the inverted-U (peak L13, $R^2 = 0.596$), with the same peak layer as the MLP probe and a trained-vs-reference gap of $\Delta = 0.103$ ($p < 0.001$).
These results indicate that our qualitative findings reflect representational structure rather than probe capacity artifacts (Figure~\ref{fig:linear_vs_mlp}).

\begin{figure}[h]
\centering
\includegraphics[width=\columnwidth]{figures/fig_linear_vs_mlp_probe.pdf}
\caption{Per-layer action probe $R^2$ for linear (dashed) and MLP (solid) probes across trained (blue) and reference (red) conditions. The inverted-U reproduces under both probe architectures; the trained--reference gap is larger for linear probes.}
\label{fig:linear_vs_mlp}
\end{figure}

\subsection{Functional Probe Comparison}
Figure~\ref{fig:probe_bars} compares functional probe $R^2$ across representation sources at each source's best layer.

\begin{figure}[h]
\centering
\includegraphics[width=\columnwidth]{figures/fig4.pdf}
\caption{Functional probe $R^2$ for action prediction (best layer per source). Trained readout ($R^2{=}0.678$, L13) and trained image\_mean ($0.655$, L8) exceed the projector baseline ($0.630$); the inverted-U extends to functional capacity.}
\label{fig:probe_bars}
\end{figure}

\subsection{Training Effect and Instruction Invariance}
Figure~\ref{fig:factorial} shows trained vs.\ reference RSA profiles and instruction invariance.
Instruction invariance holds across all conditions: maximum $|\Delta| < 0.006$, so action geometry is visually grounded and not modulated by language instructions.

\begin{figure}[h]
\centering
\includegraphics[width=\columnwidth]{figures/fig5.pdf}
\caption{\textbf{(a)}~Trained vs.\ reference image\_mean RSA profiles. Training shifts the peak from L0 ($0.612$) to L12 ($0.628$), yielding a modest best-vs-best gain ($\Delta{=}+0.017$, $3\%$). \textbf{(b)}~Instruction invariance: all conditions within $|\Delta|{<}0.006$; action geometry is visually grounded.}
\label{fig:factorial}
\end{figure}

\subsection{Additional RSA Values}
DINOv2 \texttt{patch\_mean} achieves RSA $\rho = 0.631$, between the \texttt{cls\_token} ($0.692$) and the fused DinoSigLIP ($0.634$).
VLA reference readout peaks at L1 ($\rho = 0.626$).

\subsection{Cross-Suite Validation within LIBERO: Spatial}
\label{sec:cross_benchmark_appendix}

To assess generalization beyond in-distribution tasks, we repeat the core RSA and probe analyses on LIBERO-Spatial (10 out-of-distribution tasks $\times$ 50 demonstrations each, requiring spatial reasoning).
Figure~\ref{fig:cross_benchmark} shows per-layer RSA and $R^2$ profiles for both LIBERO-Goal and LIBERO-Spatial side by side.

The inverted-U profile replicates on LIBERO-Spatial: trained image\_mean RSA peaks at L14 ($\rho = 0.691$) and $R^2$ at L13 ($0.696$), with peaks shifting ${\sim}2$ layers deeper than LIBERO-Goal (RSA L12, $R^2$ L8--L10), consistent with increased task difficulty requiring deeper processing.
Training produces a significant geometric gain (best-vs-best $\Delta = +0.023$, 95\% CI $[0.012, 0.036]$, $p < 0.001$), larger than the LIBERO-Goal effect ($\Delta = +0.017$, CI includes zero).

Two notable differences are observed.
First, the absolute RSA values are higher on Spatial ($0.691$ vs.\ $0.628$), possibly reflecting greater action diversity in spatial tasks that produces a more discriminative action RDM.
Second, RSA--$R^2$ peak convergence does not replicate: readout RSA peaks at L15 while $R^2$ peaks at L21 (6-layer gap), contrasting with the near-alignment on Goal ($|\Delta_\text{layer}| = 1$).
This suggests that peak convergence is benchmark-specific rather than a general consequence of training.

\begin{figure}[t]
\centering
\includegraphics[width=\columnwidth]{figures/fig7.pdf}
\caption{Cross-suite validation: per-layer RSA and $R^2$ profiles on LIBERO-Goal (left) vs.\ LIBERO-Spatial (right). The inverted-U replicates across suites. LIBERO-Spatial shows a significant training effect (RSA best-vs-best $\Delta = +0.023$, $p < 0.001$) and deeper peak layers (RSA L14, $R^2$ L13), but RSA--$R^2$ peak convergence does not replicate (6-layer gap).}
\label{fig:cross_benchmark}
\end{figure}

\paragraph{Cross-Suite Counterfactual Replication}
Table~\ref{tab:counterfactual_spatial} reports counterfactual instruction analysis on LIBERO-Spatial (10 tasks).
The ordering differs from LIBERO-Goal: empty $>$ shuffled $>$ wrong $>$ paraphrased (vs.\ wrong $>$ empty $>$ shuffled $>$ paraphrased on Goal).
All global effect sizes exceed $d = 0.8$ (per-task wrong minimum is $0.591$).
Spatial instructions differ only in prepositions (e.g., ``on the plate'' vs.\ ``in the drawer''); the wrong-task $d = 0.885$ suggests that instruction-dominant heads encode fine-grained spatial relations, not merely task identity.

\begin{table}[h]
\centering
\caption{Counterfactual instruction analysis on LIBERO-Spatial (10 tasks). Ordering: empty $>$ shuffled $>$ wrong $>$ paraphrased, differing from Goal. All global $d > 0.8$.}
\label{tab:counterfactual_spatial}
\small
\begin{tabular}{lccc}
\toprule
\textbf{Condition} & \textbf{Jaccard} & \textbf{Cohen's $d$} & \textbf{CV} \\
\midrule
Empty (no instruction) & 0.000 & 1.173 & 0.087 \\
Wrong (different task) & 0.174 & 0.885 & 0.272 \\
Shuffled (random order) & 1.000 & 0.979 & 0.162 \\
Paraphrased & 0.533 & 0.807 & 0.179 \\
\bottomrule
\end{tabular}
\end{table}

\paragraph{Cross-Instruction Probe on LIBERO-Spatial}
Table~\ref{tab:spatial_cross_instruction} reports cross-instruction 10-way classification accuracy on LIBERO-Spatial.
The core pattern from LIBERO-Goal replicates: image\_mean remains at chance (${\sim}10\%$) across all layers and conditions, while the trained readout token saturates to $100\%$ by L4.
The reference readout decays more steeply than on LIBERO-Goal ($75.4\%$ at L31 vs.\ $82\%$), consistent with greater task difficulty.

\begin{table}[h]
\centering
\caption{Cross-instruction 10-way classification accuracy on LIBERO-Spatial. Trained readout achieves $100\%$ from L4; image\_mean remains at chance. Reference readout decays more steeply than on LIBERO-Goal (cf.\ \S\ref{sec:openvla_routing}).}
\label{tab:spatial_cross_instruction}
\small
\begin{tabular}{llccccc}
\toprule
\textbf{Token} & \textbf{Condition} & \textbf{L1} & \textbf{L4} & \textbf{L8} & \textbf{L16} & \textbf{L31} \\
\midrule
readout & Trained & 98.9\% & 100.0\% & 100.0\% & 100.0\% & 100.0\% \\
readout & Reference & 99.5\% & 98.3\% & 85.7\% & 68.3\% & 75.4\% \\
image\_mean & Trained & 10.2\% & 10.0\% & 9.9\% & 10.0\% & 10.1\% \\
image\_mean & Reference & 10.1\% & 10.0\% & 10.0\% & 10.0\% & 10.2\% \\
\bottomrule
\end{tabular}
\end{table}

\subsection{Robustness Analyses}
\paragraph{Distance Metric Invariance}
The inverted-U profile is robust to distance metric choice: cosine action distance yields identical peak layers (L12) and Pearson correlation $r > 0.989$ with the Euclidean profile reported in the main text (Table~\ref{tab:metric_invariance}).

\begin{table}[h]
\centering
\caption{Peak layer comparison between Euclidean and cosine action distances. Profile $r$ denotes Pearson correlation across all 32 layers.}
\label{tab:metric_invariance}
\small
\begin{tabular}{lccc}
\toprule
\textbf{Condition} & \textbf{Eucl.\ Peak} & \textbf{Cos.\ Peak} & \textbf{Profile $r$} \\
\midrule
Trained image\_mean & L12 & L12 & 0.990 \\
Trained readout & L12 & L13 & 0.990 \\
Reference image\_mean & L0 & L0 & 0.995 \\
Reference readout & L1 & L1 & 0.992 \\
\bottomrule
\end{tabular}
\end{table}

\paragraph{Per-DOF Decomposition}
RSA decomposition by action dimension indicates the inverted-U is not driven by any single degree of freedom: xyz translation (peak L13, $\rho = 0.193$), rotation (peak L12, $\rho = 0.175$), and gripper (peak L13, $\rho = 0.106$) all exhibit non-monotonic profiles.
Per-DOF RSA values are lower than the 7-D RSA because single-dimension distances carry less discriminative information than the full action space.

\paragraph{Permutation Test}
Permutation testing (10{,}000 shuffles) indicates all observed RSA values are highly significant: observed $\rho$ ($0.42$--$0.63$) exceeds the null distribution maximum ($0.08$--$0.10$) by $15$--$23$ standard deviations (all $p < 0.0001$).

\paragraph{Probe Seed Sensitivity}
Action probe $R^2$ is highly stable across random seeds.
Over 5 seeds (trained image\_mean, full-layer scan), per-layer $R^2$ standard deviation is uniformly below $0.01$ (mean $0.004$), and the peak $R^2$ varies by less than $0.4\%$ ($0.655 \pm 0.002$; Table~\ref{tab:seed_sensitivity}).
Peak layer concentrates in the L8--L10 region ($4/5$ seeds); the inverted-U shape is not a seed artifact.
Seed 1 peaks at L16, but with $R^2$ difference ${<}0.002$. The inverted-U plateau (L8--L16) is sufficiently flat that argmax is sensitive to minor noise while the profile shape is invariant.
A 15-seed replication confirms stability: trained peak remains at L9 ($R^2 = 0.637 \pm 0.004$), with all trained/Llama-2-base layers showing CI width ${<}0.02$; the randomly initialized condition exhibits higher variance (std $0.185$ at peak), consistent with random initialization instability.
Fixed-layer analysis at L9 yields $R^2 = 0.637 \pm 0.008$ (trained) and $0.635 \pm 0.011$ (Llama-2 base), confirming the inverted-U pattern.
The per-seed peak method reports higher absolute values ($0.655$) because each seed selects its optimal layer independently; both methods agree on the qualitative pattern and the trained $>$ randomly initialized ordering.

\begin{table}[h]
\centering
\caption{Probe $R^2$ peak across 5 random seeds (trained image\_mean).}
\label{tab:seed_sensitivity}
\small
\begin{tabular}{lcc}
\toprule
\textbf{Seed} & \textbf{Peak Layer} & \textbf{Peak $R^2$} \\
\midrule
0 & L9 & 0.654 \\
1 & L16 & 0.656 \\
2 & L9 & 0.653 \\
3 & L10 & 0.655 \\
42 & L9 & 0.655 \\
\midrule
Mean$\pm$std & $10.6\pm2.7$ & $0.655\pm0.002$ \\
\bottomrule
\end{tabular}
\end{table}

\paragraph{Cross-architecture $\Delta R^2$ seed sensitivity.}
To verify that $\Delta R^2$ values are not sensitive to probe initialization, we retrain MLP probes with 4 different random seeds (0, 1, 2, 42) on three representative cells.
Results show high stability: TraceVLA Goal $0.227 \pm 0.010$, TraceVLA Spatial $0.136 \pm 0.002$, OpenVLA Goal $0.013 \pm 0.001$ (mean $\pm$ std across seeds).
All standard deviations are below 0.02, confirming that the exit collapse pattern ($\Delta R^2 > 0.10$ for VLM-to-VLA vs.\ $< 0.10$ for VLA-to-VLA) is robust to probe initialization.

\paragraph{Bootstrap confidence intervals for $\Delta R^2$.}
To assess uncertainty in individual $\Delta R^2$ estimates, we compute 95\% confidence intervals via 1{,}000-iteration Monte Carlo cross-validation bootstrap.
Results for 9 cells (excluding NORA-Goal due to exit-layer instability):
\begin{itemize}[nosep]
\item \textbf{VLA-to-VLA} (4 cells): CIs range from $[-0.056, 0.106]$ (widest: OpenVLA Goal) to $[0.004, 0.103]$ (OpenVLA Long).
\item \textbf{VLM-to-VLA} (5 cells): CIs range from $[0.091, 0.224]$ (TraceVLA Spatial) to $[0.283, 0.504]$ (TraceVLA Object).
\end{itemize}
The point-estimate gap ($0.097$ vs.\ $0.156$, span $= 0.059$) cleanly separates the two groups.
At the 95\% CI level, a marginal overlap of $0.016$ exists between the upper bound of OpenVLA Goal ($0.106$) and the lower bound of TraceVLA Spatial ($0.091$).
This overlap involves only 2 of 9 cells and is driven by OpenVLA Goal's unusually wide CI (due to high variance in its exit-layer probe).
All other cell pairs show non-overlapping CIs, confirming that the exit-collapse pattern is robust to measurement uncertainty.
A permutation test assesses whether the gap between VLM-to-VLA and VLA-to-VLA $\Delta R^2$ values could arise by chance.
The 9 trained-model $\Delta R^2$ values (4 VLA-to-VLA, 5 VLM-to-VLA; NORA Goal excluded due to CI instability) are randomly reassigned to groups of size 4 and 5 ($10{,}000$ permutations); the test statistic is $\min(\text{VLM-to-VLA group}) - \max(\text{VLA-to-VLA group})$.
The observed gap ($0.070$) is exceeded by $0.8\%$ of permuted distributions ($p = 0.008$), indicating that the separation is unlikely to arise by chance.
Note that the permutation shuffles transition-type labels (VLM$\to$VLA vs.\ VLA$\to$VLA), not trained/reference labels; all 9 values come from trained models.

\paragraph{Demonstration Subsample Sensitivity}
To verify that the inverted-U pattern is not an artifact of sample size, we subsample 25 from 50 demonstrations per task and recompute both RSA and probe profiles (Table~\ref{tab:demo_sensitivity}).
The RSA peak layer remains at L12 with a profile correlation of $r = 0.970$ between 50- and 25-demo conditions (high stability).
The probe peak shifts from L9 to L11 (within the L8--L16 plateau) with $r = 0.819$; $R^2$ is more sensitive to training set size than RSA.
Both metrics show reduced absolute values (RSA $-3.6\%$, $R^2$ $-7.6\%$) as expected from halved data, but the non-monotonic profile shape is preserved.

\begin{table}[h]
\centering
\caption{Inverted-U stability under demonstration subsampling (trained image\_mean).}
\label{tab:demo_sensitivity}
\small
\begin{tabular}{lccc}
\toprule
\textbf{Metric} & \textbf{50 demos} & \textbf{25 demos} & \textbf{$\Delta$} \\
\midrule
RSA peak layer & L12 & L12 & 0 \\
RSA peak $\rho$ & 0.628 & 0.606 & $-3.6\%$ \\
$R^2$ peak layer & L9 & L11 & $+2$ \\
$R^2$ peak value & 0.659 & 0.609 & $-7.6\%$ \\
RSA profile $r$ & --- & 0.970 & --- \\
$R^2$ profile $r$ & --- & 0.819 & --- \\
\bottomrule
\end{tabular}
\end{table}

\paragraph{Label Permutation Sanity Check}
\label{sec:label_permutation}

To confirm that probe performance reflects genuine action-relevant structure rather than overfitting, we repeat the per-layer action probe with randomly shuffled labels (Figure~\ref{fig:label_permutation}).
Under permuted labels, Spearman $\rho$ drops from $0.955$ to $0.001$ (trained) and from $0.933$ to $0.000$ (reference), with all 33 layers producing $|\rho| < 0.013$.
This establishes that the high probe accuracy reported throughout the paper is attributable to learned action-relevant representations, not probe capacity artifacts.

\begin{figure}[h]
\centering
\includegraphics[width=0.55\columnwidth]{figures/fig_label_permutation.pdf}
\caption{\textbf{Label permutation control.} Action-probe Spearman $\rho$ under real versus shuffled labels for OpenVLA trained and reference conditions. Shuffled-label $\rho$ is indistinguishable from zero, confirming that probe performance reflects genuine representational structure. (Evaluated with random frame splitting; demo-level split yields lower absolute values but identical null-vs-real separation; see \S\ref{sec:cross_arch_analyses}.)}
\label{fig:label_permutation}
\end{figure}

\paragraph{Statistical Power}
Post-hoc power analysis ($N = 10$ tasks, $\alpha = 0.05$, target power $= 0.80$) indicates minimum detectable effect sizes of $\Delta_{\text{RSA}} = 0.030$ and $\Delta_{R^2} = 0.026$.
The observed RSA training effect ($\Delta = 0.017$, Cohen's $d_z = 0.56$) achieves power $= 0.356$, confirming the analysis is underpowered for this effect size; the non-significant RSA result (\S\ref{sec:routing_synthesis}) should not be interpreted as evidence of absence.
The observed $R^2$ readout effect ($\Delta = 0.069$, $d_z = 2.68$) achieves power $= 1.000$, confirming adequate sensitivity for the primary functional metric.

\paragraph{$\Delta R^2$ Threshold Robustness}
\label{sec:delta_r2_threshold}

The main text classifies $\Delta R^2 > 0.10$ as substantial exit collapse.
Table~\ref{tab:delta_r2_threshold} shows that any threshold in the interval $(0.097, 0.156)$ yields identical three-way separation: all VLM-to-VLA trained cells above threshold, all OpenVLA trained and all reference cells below.
This interval is bounded by the largest non-VLM-to-VLA value ($0.097$, TraceVLA Long reference) and the smallest VLM-to-VLA trained value ($0.156$, TraceVLA Spatial).
The $0.10$ threshold was not selected to maximize separation; rather, it represents a natural effect-size boundary ($\geq 10\%$ absolute change in explained variance) that happens to fall within this gap.

\begin{table}[h]
\centering
\caption{Sensitivity of group separation to $\Delta R^2$ threshold choice. Clean separation (all VLM-to-VLA trained above, all others below) holds for any threshold in $(0.097, 0.156)$; the 0.10 threshold is not boundary-sensitive.}
\label{tab:delta_r2_threshold}
\small
\begin{tabular}{ccccc}
\toprule
\textbf{Threshold} & \textbf{VLM$\to$VLA Tr.} & \textbf{OpenVLA Tr.} & \textbf{Reference} & \textbf{Clean} \\
& \textbf{($N\!/6$)} & \textbf{($N\!/4$)} & \textbf{($N\!/10$)} & \textbf{sep.} \\
\midrule
0.05 & 6/6 & 1/4 & 6/10 & No \\
0.08 & 6/6 & 1/4 & 3/10 & No \\
\textbf{0.10} & \textbf{6/6} & \textbf{0/4} & \textbf{0/10} & \textbf{Yes} \\
\textbf{0.15} & \textbf{6/6} & \textbf{0/4} & \textbf{0/10} & \textbf{Yes} \\
0.20 & 5/6 & 0/4 & 0/10 & No \\
\bottomrule
\end{tabular}
\end{table}

\subsection{Detailed Training Effects}
Table~\ref{tab:training_effects} reports bootstrap-based statistical comparisons of training effects on RSA and probe $R^2$.

\begin{table}[h]
\centering
\caption{Training effects on action-relevant structure (LIBERO-Goal; bootstrap, 10{,}000 resamples). RSA best-vs-best yields a modest geometric gain ($\Delta{=}+0.017$, $3\%$; LIBERO-Spatial: $+0.023$, $p{<}0.001$), while $R^2$ shows substantially larger functional gains.}
\label{tab:training_effects}
\small
\begin{tabular}{lccc}
\toprule
\textbf{Comparison} & \textbf{$\Delta$} & \textbf{95\% CI} & \textbf{$p$} \\
\midrule
\multicolumn{4}{l}{\emph{RSA ($\rho$)}} \\
Best-vs-best image\_mean & $+$0.017 & incl.\ 0 & n.s. \\
\quad (trained L12 vs.\ reference L0) & & & \\
Same-layer L12 image\_mean & $+$0.079 & excl.\ 0 & $<$0.0001 \\
\quad (trained vs.\ reference) & & & \\
Instruction effect (max $|\Delta|$) & 0.006 & --- & n.s. \\
\midrule
\multicolumn{4}{l}{\emph{Probe $R^2$}} \\
Best-vs-best image\_mean & $+$0.015 & excl.\ 0 & 0.006 \\
\quad (trained L8 vs.\ reference L0) & & & \\
Best-vs-best readout & $+$0.069 & excl.\ 0 & $<$0.0001 \\
\quad (trained L13 vs.\ reference L32) & & & \\
\bottomrule
\end{tabular}
\end{table}

\subsection{Training Effect Size Summary}
Table~\ref{tab:training_effect_summary} consolidates training effects across all metrics.
Training effects on image\_mean metrics are modest (RSA $\Delta = +0.017$, $R^2$ $\Delta = +0.015$), but effects on readout metrics are substantial (instruction attention $+18$pp, attention $+46$pp, $R^2$ $\Delta = +0.069$).
This asymmetry indicates that training primarily changes how the readout token integrates information rather than altering the visual representations carried by image tokens.

\begin{table}[h]
\centering
\caption{Training effect sizes across metrics. Image\_mean effects are modest; readout effects are large. Training targets readout-token information integration.}
\label{tab:training_effect_summary}
\small
\begin{tabular}{lccc}
\toprule
\textbf{Metric} & \textbf{Trained} & \textbf{Reference} & \textbf{$\Delta$} \\
\midrule
RSA peak (image\_mean) & 0.628 (L12) & 0.612 (L0) & $+$0.017 \\
$R^2$ peak (image\_mean) & 0.655 (L8) & 0.640 (L0) & $+$0.015 \\
$R^2$ peak (readout) & 0.678 (L13) & 0.609 (L32) & $+$0.069 \\
Instruction attention L31 & 100\% & 82\% & $+$18pp \\
Instruction attention L1 & 58\% & 12\% & $+$46pp \\
Progress classification & 86.3\% (L20) & 85.3\% (L4) & $+$1pp \\
\bottomrule
\end{tabular}
\end{table}

\subsection{Instruction Classification Probes}
\paragraph{Cross-Instruction Classification: Dense Layer Profile}
\label{sec:dense_routing}

Table~\ref{tab:dense_routing} shows the dense per-layer profile of cross-instruction classification (10-way) at the readout token for layers L0--L8.
The step-function onset between L0 ($10\%$, chance) and L1 (${\geq}99.6\%$) reflects the high discriminability of instruction embeddings from the frozen tokenizer: even incidental attention (${\sim}12\%$) suffices for linear classification at L1.
The VLA-trained model sustains $100\%$ through L31; Llama-2 base reaches $100\%$ by L5 and maintains it. The OXE-pretrained reference and random initialization both decay to ${\sim}82\%$ by L31, indicating that deep-layer instruction maintenance requires either VLA training or structured language-model pretraining on text-only data.

\begin{table}[h]
\centering
\caption{Cross-instruction 10-way classification accuracy at readout token, L0--L8. Values reflect dense classification evaluation; the main text (\S\ref{sec:openvla_routing}) reports sparse single-split results, yielding minor differences (e.g., L8: $98.2\%$ vs.\ $98.6\%$).}
\label{tab:dense_routing}
\small
\begin{tabular}{lccc}
\toprule
\textbf{Layer} & \textbf{Trained} & \textbf{Reference} & \textbf{Image\_mean (trained)} \\
\midrule
L0 & 10.0\% & 10.0\% & 10.0\% \\
L1 & 100.0\% & 100.0\% & 10.0\% \\
L2 & 100.0\% & 99.6\% & 10.0\% \\
L3 & 100.0\% & 99.6\% & 10.0\% \\
L4 & 100.0\% & 99.1\% & 10.0\% \\
L5 & 100.0\% & 98.2\% & 10.0\% \\
L6 & 100.0\% & 97.3\% & 10.0\% \\
L7 & 100.0\% & 96.4\% & 10.0\% \\
L8 & 100.0\% & 98.6\% & 10.0\% \\
\bottomrule
\end{tabular}
\end{table}

\paragraph{Cross-Instruction LOOCV Validation}
\label{sec:loocv}

Leave-one-out cross-validation over 10 folds (each fold holds out one task's instructions) validates that the cross-instruction probe result is not driven by memorization.
Trained readout achieves $100\% \pm 0.00$ across all 10 folds at layers L1--L31 (perfect generalization).

\paragraph{QA Probe: Progress and Subtask Classification}
To test whether LLM layers encode task-relevant semantic information beyond instruction identity, we train logistic regression probes (scikit-learn, $C = 1.0$, max\_iter $= 1{,}000$) to classify task progress (4-class: each trajectory split into 4 equal temporal segments) and active subtask (binary: first vs.\ second half) from frozen representations, with an 80/20 random train/test split (Table~\ref{tab:qa_probe}).

\begin{table}[h]
\centering
\caption{QA probe accuracy (progress + subtask classification, best layer).}
\label{tab:qa_probe}
\small
\begin{tabular}{lcc}
\toprule
\textbf{Source} & \textbf{Best Layer} & \textbf{Accuracy} \\
\midrule
Trained readout & L20 & 86.3\% \\
Reference readout & L20 & 85.3\% \\
DINOv2 cls\_token & --- & 82.7\% \\
\bottomrule
\end{tabular}
\end{table}

\paragraph{Cross-Instruction Probe Profiles}
\begin{figure}[h]
\centering
\includegraphics[width=\columnwidth]{figures/fig_probe_conditions.pdf}
\caption{Cross-instruction 10-way classification accuracy at the readout token across all 32 LLM layers. VLA-trained achieves $99.6\%$ at L1, saturating to $100\%$. Llama-2 base starts at $80.2\%$, reaching $100\%$ by L5. Reference starts at $100\%$ but decays to $82\%$ by L31.}
\end{figure}

\paragraph{Cross-Architecture Probe and Attention Profile}
Table~\ref{tab:cross_arch_probe} reports cross-instruction probe accuracy and L1 instruction-dominant head counts across six conditions spanning two architectures.
The probe-attention dissociation at L1 is notable: TraceVLA training slightly \emph{reduces} instruction-dominant heads ($32 \to 27$--$30/32$) while \emph{decreasing} probe accuracy ($98.1\% \to 83.6\%$), consistent with representational reorganization rather than linear separability.

\begin{table}[h]
\centering
\caption{Cross-instruction 10-way probe accuracy and L1 instruction-dominant heads across six conditions.
Probe accuracy and attention head count are dissociated at L1: TraceVLA training slightly reduces heads ($32 \to 27$--$30/32$) while decreasing probe accuracy ($98.1\% \to 83.6\%$), paralleling the information-computation dissociation observed in OpenVLA.
Image\_mean remains at chance ($10\%$) in all conditions (omitted).}
\label{tab:cross_arch_probe}
\small
\begin{tabular}{lccccc}
\toprule
\textbf{Condition} & \textbf{L1 heads} & \textbf{L1} & \textbf{L2+} & \textbf{L31} & \textbf{Mode} \\
\midrule
\multicolumn{6}{l}{\emph{OpenVLA (Llama-2 7B)}} \\
\quad Random init & 0/32 & 100.0\% & decays & 82.0\% & incid.+forget \\
\quad Llama-2 base & 9/32 & 80.2\% & L5: 100\% & 100.0\% & gradual+stable \\
\quad VLA trained & 24/32 & 99.6\% & 100\% & 100.0\% & immed.+stable \\
\quad LIBERO-finetuned & 22/32 & 99.7\% & 100\% & 100.0\% & immed.+stable \\
\midrule
\multicolumn{6}{l}{\emph{TraceVLA (Phi-3-V 4B)}} \\
\quad Phi-3-Vision & 32/32 & 98.1\% & 100\% & 100.0\% & inherited+stable \\
\quad VLA trained & 27--30/32 & 83.6\% & 100\% & 100.0\% & reorg.+stable \\
\bottomrule
\end{tabular}
\end{table}

\begin{figure}[h]
\centering
\includegraphics[width=\columnwidth]{figures/fig_cross_architecture_attention.pdf}
\caption{Instruction-dominant heads across layers for two architectures (OpenVLA / TraceVLA), each comparing pre-VLA and VLA-trained conditions. Despite opposite starting points (OpenVLA/Llama-2 begins with sparse attention; TraceVLA/Phi-3-V begins with distributed attention), VLA training drives both to entry-exit concentration (L1 saturation, mid-layer suppression, L29--L31 re-engagement).}
\end{figure}

\paragraph{Cross-Architecture L1 Instruction Dominance}
Figure~\ref{fig:cross_arch_l1_bar} complements the main-text summary (Figure~\ref{fig:cross_arch_l1}) by showing the fraction of instruction-dominant heads at L1 across all three architectures in a single comparison. The trained/reference fractions are comparable within each architecture (OpenVLA: 23/32 vs.\ 26/32; TraceVLA: 29/32 vs.\ 32/32; NORA: 13/16 vs.\ 14/16), consistent with L1 instruction attention being largely inherited from the pretrained backbone rather than acquired during VLA fine-tuning.

\begin{figure}[h]
\centering
\includegraphics[width=0.65\columnwidth]{figures/fig_cross_arch_l1.pdf}
\caption{\textbf{L1 instruction-dominant head fractions across three architectures.} Trained and reference conditions show comparable fractions in all three VLAs, indicating that L1 instruction attention largely pre-exists in the pretrained backbone.}
\label{fig:cross_arch_l1_bar}
\end{figure}

\paragraph{L1 Attention Pattern Analysis}
Figure~\ref{fig:attention_appendix} shows the per-head L1 attention distribution at the readout token.
The three-level gradient (randomly initialized 0/32, Llama-2 base 9/32, VLA-trained 24/32 instruction-dominant heads) shows that language pretraining seeds partial instruction attention and VLA fine-tuning shifts it further toward instruction dominance.

\begin{figure}[h]
\centering
\includegraphics[width=\columnwidth]{figures/fig_l1_stacked_bar.pdf}
\caption{L1 attention from the readout token across three conditions, per head (sorted by instruction attention fraction). VLA-trained: 24/32 instruction-dominant heads ($57.7\%$ text attention). Llama-2 base: 9/32 heads ($36.5\%$). Randomly initialized: 0/32 heads ($11.7\%$). See Table~\ref{tab:multilayer_attention} for quantitative comparison across layers.}
\label{fig:attention_appendix}
\end{figure}

\begin{figure}[h]
\centering
\includegraphics[width=0.6\columnwidth]{figures/fig_l1_attention.pdf}
\caption{L1 instruction-dominant heads (trained vs.\ reference). VLMs saturate; training does not increase counts.}
\label{fig:cross_arch_l1}
\end{figure}

\paragraph{Multi-Layer Attention Patterns}
Table~\ref{tab:multilayer_attention} shows a 4-layer subset of the full 32-layer attention profile.
The full profile (Figure~\ref{fig:synthesis_panels}) reveals three qualitatively distinct patterns: the randomly initialized model maintains 0 instruction-dominant heads at all 32 layers; Llama-2 base shows a mid-depth peak (L9--L16: 15--22 heads), consistent with typical autoregressive LM processing depth; VLA fine-tuning produces an entry-exit pattern (L0--L1: 21--24 heads, L31: 18 heads, intermediate layers 0--11).

% Figure promoted to main text (Figure~\ref{fig:multilayer_attention} in \S\ref{sec:synthesis}).

\begin{table}[h]
\centering
\caption{Instruction-dominant heads and text attention fraction at four sampled layers (subset of the 32-layer profile in Figure~\ref{fig:synthesis_panels}). VLA-trained shows entry-exit concentration (L1, L31); Llama-2 base peaks at mid-depth (L16); randomly initialized shows no instruction attention. Head counts from the 32-layer bulk extraction may differ slightly from dedicated per-layer analyses (\S\ref{sec:openvla_routing}) due to threshold boundary effects.}
\label{tab:multilayer_attention}
\small
\begin{tabular}{lcccccc}
\toprule
 & \multicolumn{2}{c}{\textbf{VLA Trained}} & \multicolumn{2}{c}{\textbf{Llama-2 base}} & \multicolumn{2}{c}{\textbf{Rand.\ Init}} \\
\cmidrule(lr){2-3} \cmidrule(lr){4-5} \cmidrule(lr){6-7}
\textbf{Layer} & \textbf{Heads} & \textbf{Text\%} & \textbf{Heads} & \textbf{Text\%} & \textbf{Heads} & \textbf{Text\%} \\
\midrule
L1 & \textbf{24/32} & 57.7 & 9/32 & 36.5 & 0/32 & 11.7 \\
L8 & 4/32 & 46.8 & 8/32 & 48.8 & 0/32 & 7.9 \\
L16 & 10/32 & 48.4 & \textbf{25/32} & 60.2 & 0/32 & 9.2 \\
L31 & \textbf{18/32} & 56.3 & 8/32 & 34.9 & 0/32 & 7.6 \\
\bottomrule
\end{tabular}
\end{table}

\paragraph{Threshold Sensitivity Analysis}
\label{sec:threshold_sensitivity}

The main text defines instruction-dominant heads as those allocating $>50\%$ attention to instruction tokens.
Table~\ref{tab:threshold_sensitivity} shows that the three-condition gradient (randomly initialized $<$ Llama-2 base $<$ VLA-trained) holds across all thresholds from 30\% to 70\%; the 50\% threshold is not a boundary-sensitive choice.

\begin{table}[h]
\centering
\caption{Number of instruction-dominant heads at L1 under varying attention thresholds. The three-condition ordering is stable across all thresholds.}
\label{tab:threshold_sensitivity}
\small
\begin{tabular}{lccc}
\toprule
\textbf{Threshold} & \textbf{VLA Trained} & \textbf{Llama-2 base} & \textbf{Rand.\ Init} \\
\midrule
30\% & 32 & 13 & 0 \\
40\% & 31 & 10 & 0 \\
50\% & 24 & 9 & 0 \\
60\% & 9 & 4 & 0 \\
70\% & 5 & 3 & 0 \\
\bottomrule
\end{tabular}
\end{table}

\subsection{L1 Attention Decomposition by Head Group}
\label{sec:attention_decomp}

Table~\ref{tab:attention_decomp} decomposes L1 attention mass at the readout token into three token categories for instruction-dominant and control head groups.

\begin{table}[h]
\centering
\caption{Attention mass allocation at Layer 1, averaged across 5 samples. Instruction-dominant heads (defined as $>50\%$ attention to instruction tokens) allocate 17.5 percentage points more attention to instruction tokens than control heads, with corresponding reductions in vision ($-8.3$pp) and BOS/sink ($-9.2$pp) attention.}
\label{tab:attention_decomp}
\small
\begin{tabular}{lcccc}
\toprule
\textbf{Head Group} & \textbf{$N$} & \textbf{Instruction} & \textbf{Vision} & \textbf{BOS/Other} \\
\midrule
Instruction-dominant & 24 & 62.4\% & 18.8\% & 18.8\% \\
Control & 8 & 44.9\% & 27.1\% & 28.0\% \\
\midrule
$\Delta$ & --- & $+$17.5pp & $-$8.3pp & $-$9.2pp \\
\bottomrule
\end{tabular}
\end{table}

Head-level instruction attention entropy reaches $99.5\%$ of the uniform maximum; attention is broadly distributed rather than concentrated in a few heads.

\paragraph{Per-token attention normalization.}
The $57.7\%$ aggregate text attention in instruction-dominant heads partly reflects token-count imbalance (22 instruction vs.\ 256 image tokens).
Per-token normalization reveals a $38.2\times$ ratio: each instruction token receives $38.2\times$ the attention of each image token in trained instruction-dominant heads, compared to only $1.5\times$ in randomly initialized heads (a $25.5\times$ increase from training).
Instruction-dominant heads also show lower attention entropy ($0.465$) than non-dominant heads ($0.569$), indicating more concentrated attention distributions.

\subsection{Per-Dimension Action MAE Under Head Ablation}
Table~\ref{tab:per_dim_mae} extends Table~\ref{tab:action_shift_perdim} with full per-dimension control-head data and instruction-to-control MAE ratios.

\begin{table}[h]
\centering
\caption{Mean absolute error per action dimension under L1 head ablation. Translation dimensions ($dx$/$dy$/$dz$) show 3--5$\times$ larger instruction-to-control ratio than rotation dimensions (1.9--2.6$\times$); instruction-dominant heads have disproportionate influence on positional control. Gripper (2.07$\times$) falls between the two groups.}
\label{tab:per_dim_mae}
\small
\begin{tabular}{lccc}
\toprule
\textbf{Dimension} & \textbf{Instr.\ 24h MAE} & \textbf{Control 8h MAE} & \textbf{Ratio} \\
\midrule
$dx$ & 0.00189 & 0.00037 & 5.09$\times$ \\
$dy$ & 0.00531 & 0.00125 & 4.26$\times$ \\
$dz$ & 0.00762 & 0.00243 & 3.14$\times$ \\
$roll$ & 0.02059 & 0.00785 & 2.62$\times$ \\
$pitch$ & 0.03071 & 0.01387 & 2.21$\times$ \\
$yaw$ & 0.03388 & 0.01783 & 1.90$\times$ \\
$gripper$ & 0.29072 & 0.14072 & 2.07$\times$ \\
\bottomrule
\end{tabular}
\end{table}

\subsection{Per-Task $R^2$ Distribution Under L1 Ablation}
Table~\ref{tab:per_task_r2} reports per-task action probe $R^2$ at L1 under baseline and 24-head ablation conditions (MLP probe, 5 demonstrations per task).

\begin{table}[h]
\centering
\caption{Per-task action probe $R^2$ at L1 (MLP probe, 5 demos/task). All 10 tasks show positive ablation effect ($\Delta > 0$): instruction-head ablation uniformly increases action decodability at the entry layer. Mean $\Delta R^2 = +0.107$; the 50-demo setting (\S\ref{sec:openvla_routing}) yields a consistent $+0.059$.}
\label{tab:per_task_r2}
\small
\begin{tabular}{lccc}
\toprule
\textbf{Task} & \textbf{Baseline $R^2$} & \textbf{Ablated $R^2$} & \textbf{$\Delta R^2$} \\
\midrule
open middle drawer & 0.197 & 0.221 & $+$0.024 \\
open top drawer + bowl & 0.265 & 0.340 & $+$0.075 \\
push plate to stove & 0.190 & 0.197 & $+$0.007 \\
put bowl on plate & 0.375 & 0.474 & $+$0.099 \\
put bowl on stove & 0.235 & 0.395 & $+$0.160 \\
put bowl on cabinet & 0.256 & 0.457 & $+$0.201 \\
put cream cheese in bowl & 0.186 & 0.345 & $+$0.159 \\
put wine on rack & 0.189 & 0.424 & $+$0.235 \\
put wine on cabinet & 0.279 & 0.353 & $+$0.074 \\
turn on stove & 0.385 & 0.421 & $+$0.036 \\
\midrule
Mean $\pm$ std & $0.256 \pm 0.070$ & $0.363 \pm 0.088$ & $+$0.107 \\
\bottomrule
\end{tabular}
\end{table}

\subsection{Task Complexity and Attention Dependence}
\label{sec:task_complexity}

We classify tasks as \emph{instruction-discriminative} when multiple tasks share the same manipulated object(s), requiring instruction content to disambiguate the target action (e.g., ``put bowl on plate'' vs.\ ``put bowl on stove''); remaining tasks involve unique object-action pairings recoverable from visual context alone.
Table~\ref{tab:task_complexity} reports per-task $|\Delta R^2|$ under both groupings.

\begin{table}[h]
\centering
\caption{Task-level L1 ablation effects ($|\Delta R^2|$) grouped by instruction dependence. \emph{Primary} (conservative): cream cheese in bowl classified as visual-sufficient ($2.5\times$, $p = 0.033$). \emph{Sensitivity} (shared-object criterion): cream cheese reclassified as instruction-discriminative ($6.5\times$, $p = 0.008$).}
\label{tab:task_complexity}
\small
\begin{tabular}{lcc}
\toprule
\textbf{Task} & \textbf{$|\Delta R^2|$} & \textbf{Group} \\
\midrule
put wine on rack & 0.235 & Instr-disc \\
put bowl on cabinet & 0.201 & Instr-disc \\
put bowl on stove & 0.160 & Instr-disc \\
put bowl on plate & 0.099 & Instr-disc \\
open top drawer + bowl & 0.075 & Instr-disc \\
put wine on cabinet & 0.074 & Instr-disc \\
\midrule
put cream cheese in bowl$^\dagger$ & 0.159 & Vis-suff \\
turn on stove & 0.036 & Vis-suff \\
open middle drawer & 0.024 & Vis-suff \\
push plate to stove & 0.007 & Vis-suff \\
\midrule
\multicolumn{3}{l}{\textbf{Primary}: Instr-disc mean $0.141$; Vis-suff mean $0.057$} \\
\multicolumn{3}{l}{Mann-Whitney $U = 21.0$, $p = 0.033$, ratio $= 2.49\times$} \\
\multicolumn{3}{l}{\textbf{Sensitivity}$^\dagger$: cream cheese as Instr-disc: $p = 0.008$, $6.5\times$} \\
\bottomrule
\end{tabular}
\end{table}

\paragraph{Leave-one-out sensitivity.}
All leave-one-out $p$-values remain $\leq 0.083$ (marginally significant), with no single task removal reversing the effect direction.
Reclassifying the boundary case (cream cheese in bowl) to the instruction-discriminative group strengthens the result ($p = 0.008$).

Bootstrap 95\% CI on the complex/simple ratio: $[1.03, 11.12]$ (10{,}000 resamples). The wide interval reflects low statistical power from $N=10$ tasks split into groups of 6 and 4; the point estimate ($2.5\times$) persists under leave-one-out analysis (all $p \leq 0.083$, marginally significant; \S\ref{sec:task_complexity}).

\subsection{Specificity Control: Matched-Size Head Ablation}
To test whether ablation effects are specific to instruction-dominant heads, we ablate four matched-size subsets of 24 heads at L1: instruction-dominant (highest instruction attention fraction), image-dominant (highest image attention fraction), random (5 seeds), and bottom-24 (lowest instruction attention fraction; overlaps 16/24 heads with the instruction-dominant group).

\begin{table}[h]
\centering
\caption{Matched-size ablation controls: trajectory disruption ($d$) from ablating 24-head subsets at L1, ranked by instruction attention fraction. Instruction-dominant heads produce the largest effect, but the modest gap ($8\%$ over random) indicates distributed, attention-graded redistribution rather than a dedicated subset.}
\label{tab:specificity_control}
\small
\begin{tabular}{lcc}
\toprule
\textbf{Condition (24 heads)} & \textbf{Cohen's $d$} & \textbf{95\% CI} \\
\midrule
Instruction-dominant & 0.650 & (reference) \\
Random (5 seeds) & $0.601 \pm 0.021$ & $[0.568, 0.622]$ \\
Image-dominant & 0.580 & $[0.547, 0.613]$ \\
Bottom-24 & 0.565 & --- \\
\bottomrule
\end{tabular}
\end{table}

All conditions yield comparable ground-truth MSE ($0.487$--$0.504$ vs.\ baseline $0.497$), consistent with direction-selective rather than precision-degrading effects.

\subsection{Paired Effect Sizes}
Main-text Cohen's $d$ values use the standard (unpaired) formula for token-level comparisons where observations are independent tokens.
For within-subject effects where the same tasks or demonstrations are measured under two conditions, we additionally report paired $d_z = \bar{\Delta} / s_\Delta$, which accounts for within-subject correlation and yields larger effect sizes when correlation is high (Table~\ref{tab:paired_dz}).

\begin{table}[h]
\centering
\caption{Paired effect sizes ($d_z$) for within-subject comparisons. Standard (unpaired) $d$ reported where available. Paired $d_z$ exceeds standard $d$ when within-subject correlation is high.}
\label{tab:paired_dz}
\small
\begin{tabular}{lccc}
\toprule
\textbf{Comparison} & \textbf{Standard $d$} & \textbf{Paired $d_z$} & \textbf{$N$} \\
\midrule
L1 ablation: 24 dominant vs.\ baseline & 0.65 & 1.17 & 10 tasks \\
Random-24 ablation vs.\ baseline & --- & $1.84 \pm 0.33$ & 30 subsets \\
Dose-response: 32 vs.\ 4 heads (MSE) & --- & 1.79 & 10 tasks \\
RSA: trained vs.\ reference (best-vs-best) & --- & 0.55 & 10 tasks \\
Cross-arch: OpenVLA vs.\ TraceVLA peak & --- & 0.48 & 15 seeds \\
\bottomrule
\end{tabular}
\end{table}

\subsection{Cross-Architecture Action Decodability: $2 \times 2$ Design}
\label{sec:cross_arch_r2}

Table~\ref{tab:2x2_r2} reports per-layer action prediction $R^2$ across two architectures (OpenVLA, TraceVLA) $\times$ two training conditions (VLA-trained, reference base LLM), providing a $2 \times 2$ factorial view of how per-layer action decodability differs between trained and reference conditions.

\paragraph{Convergent reorganization from divergent starting points.}
Reference profiles differ qualitatively: OpenVLA (Llama-2) peaks at L0 ($R^2 = 0.629$) then collapses to near-zero at mid-layers (L20: $0.086$, std $= 0.311$), while Phi-3-V peaks at L6 ($0.578 \pm 0.018$) with gradual monotonic decline (late-layer std $= 0.068$).
After VLA training, both architectures reach nearly identical mid-layer peaks (OpenVLA at L9, $R^2 = 0.637 \pm 0.008$; TraceVLA at L11, $0.634 \pm 0.009$) with $\Delta = 0.003$ and overlapping CIs.
Deep-layer variance drops sharply with training: OpenVLA std $0.277 \to 0.040$; Phi-3-V std $0.068 \to 0.011$.
Peak $R^2$ values are comparable before and after training (reference peaks reach $91$--$99\%$ of trained), but peak locations shift deeper (L0$\to$L9; L6$\to$L11) and profile shapes transform from architecture-specific patterns to a shared mid-layer plateau.

\paragraph{Summary.}
Table~\ref{tab:2x2_r2_summary} condenses the key contrasts.
The L0 values are identical between trained and reference within each architecture (shared frozen vision encoder + projector), confirming that training effects arise from LLM weight changes.
The largest training gains appear at mid-to-deep layers (L15--L25 average: OpenVLA $+0.313$, Phi-3-V $+0.103$), precisely where reference profiles diverge most.

\begin{table}[h]
\centering
\caption{$R^2$ profile summary across architectures and training conditions.}
\label{tab:2x2_r2_summary}
\begin{tabular}{lcccc}
\toprule
 & \multicolumn{2}{c}{\textbf{OpenVLA (Llama-2-7B)}} & \multicolumn{2}{c}{\textbf{TraceVLA (Phi-3-mini)}} \\
\cmidrule(lr){2-3} \cmidrule(lr){4-5}
 & \textbf{Reference} & \textbf{Trained} & \textbf{Reference} & \textbf{Trained} \\
\midrule
Peak Layer       & L0          & L9          & L6          & L11         \\
Peak $R^2$       & 0.629       & 0.637       & 0.578       & 0.634       \\
Profile Shape    & Collapse    & Flat plateau & Ramp $\to$ decline & Ramp $\to$ plateau \\
Late-layer std   & 0.277       & 0.040       & 0.068       & 0.011       \\
$\Delta_{\text{train}}$ (peak) & \multicolumn{2}{c}{$+0.008$ (L0$\to$L9)} & \multicolumn{2}{c}{$+0.056$ (L6$\to$L11)} \\
$\Delta_{\text{train}}$ (L15--L25 avg)$^\dagger$ & \multicolumn{2}{c}{$+0.313$} & \multicolumn{2}{c}{$+0.103$} \\
\bottomrule
\end{tabular}
\vspace{2pt}
{\footnotesize $^\dagger$Computed from full-precision per-seed values before rounding; displayed per-layer values (Table~\ref{tab:2x2_r2}) may yield slightly different averages due to rounding. Late-layer std computed over L20--L31.}
\end{table}

\begin{table}[h]
\centering
\caption{Per-layer action prediction $R^2$: training effect across architectures (mean $\pm$ std over 15 seeds). Peak values per condition in \textbf{bold}.}
\label{tab:2x2_r2}
\resizebox{\textwidth}{!}{%
\begin{tabular}{lcccc}
\toprule
\textbf{Layer} & \multicolumn{2}{c}{\textbf{OpenVLA (Llama-2-7B)}} & \multicolumn{2}{c}{\textbf{TraceVLA (Phi-3-mini)}} \\
\cmidrule(lr){2-3} \cmidrule(lr){4-5}
 & \textbf{Reference} & \textbf{Trained} & \textbf{Reference} & \textbf{Trained} \\
\midrule
L0  & $\mathbf{0.629 \pm 0.016}$ & $0.629 \pm 0.016$ & $-0.330 \pm 0.003$ & $-0.330 \pm 0.002$ \\
L1  & $0.539 \pm 0.023$ & $0.625 \pm 0.017$ & $0.429 \pm 0.034$ & $0.447 \pm 0.013$ \\
L2  & $0.517 \pm 0.047$ & $0.622 \pm 0.010$ & $0.458 \pm 0.023$ & $0.522 \pm 0.014$ \\
L3  & $0.448 \pm 0.216$ & $0.624 \pm 0.014$ & $0.498 \pm 0.036$ & $0.579 \pm 0.011$ \\
L4  & $0.496 \pm 0.057$ & $0.628 \pm 0.014$ & $0.572 \pm 0.021$ & $0.592 \pm 0.009$ \\
L5  & $0.475 \pm 0.101$ & $0.627 \pm 0.019$ & $0.573 \pm 0.020$ & $0.602 \pm 0.010$ \\
L6  & $0.422 \pm 0.119$ & $0.631 \pm 0.012$ & $\mathbf{0.578 \pm 0.018}$ & $0.609 \pm 0.012$ \\
L7  & $0.459 \pm 0.093$ & $0.632 \pm 0.012$ & $0.569 \pm 0.023$ & $0.616 \pm 0.014$ \\
L8  & $0.347 \pm 0.283$ & $0.635 \pm 0.013$ & $0.560 \pm 0.025$ & $0.615 \pm 0.010$ \\
L9  & $0.445 \pm 0.127$ & $\mathbf{0.637 \pm 0.008}$ & $0.564 \pm 0.021$ & $0.617 \pm 0.013$ \\
L10 & $0.353 \pm 0.224$ & $0.624 \pm 0.014$ & $0.551 \pm 0.027$ & $0.625 \pm 0.007$ \\
L11 & $0.356 \pm 0.122$ & $0.623 \pm 0.016$ & $0.561 \pm 0.025$ & $\mathbf{0.634 \pm 0.009}$ \\
L12 & $0.409 \pm 0.203$ & $0.626 \pm 0.011$ & $0.557 \pm 0.026$ & $0.632 \pm 0.005$ \\
L13 & $0.341 \pm 0.227$ & $0.624 \pm 0.013$ & $0.547 \pm 0.028$ & $0.633 \pm 0.008$ \\
L14 & $0.374 \pm 0.211$ & $0.623 \pm 0.015$ & $0.550 \pm 0.023$ & $0.631 \pm 0.009$ \\
L15 & $0.297 \pm 0.225$ & $0.619 \pm 0.012$ & $0.550 \pm 0.022$ & $0.633 \pm 0.010$ \\
L16 & $0.323 \pm 0.226$ & $0.619 \pm 0.011$ & $0.534 \pm 0.018$ & $0.631 \pm 0.010$ \\
L17 & $0.395 \pm 0.109$ & $0.609 \pm 0.023$ & $0.519 \pm 0.024$ & $0.633 \pm 0.007$ \\
L18 & $0.214 \pm 0.307$ & $0.613 \pm 0.017$ & $0.502 \pm 0.047$ & $0.629 \pm 0.012$ \\
L19 & $0.241 \pm 0.258$ & $0.613 \pm 0.013$ & $0.499 \pm 0.058$ & $0.627 \pm 0.008$ \\
L20 & $0.086 \pm 0.311$ & $0.611 \pm 0.014$ & $0.482 \pm 0.048$ & $0.628 \pm 0.011$ \\
L21 & $0.330 \pm 0.216$ & $0.589 \pm 0.023$ & $0.476 \pm 0.036$ & $0.624 \pm 0.007$ \\
L22 & $0.207 \pm 0.283$ & $0.586 \pm 0.013$ & $0.408 \pm 0.140$ & $0.616 \pm 0.008$ \\
L23 & $0.264 \pm 0.227$ & $0.578 \pm 0.028$ & $0.455 \pm 0.074$ & $0.613 \pm 0.007$ \\
L24 & $0.173 \pm 0.275$ & $0.577 \pm 0.013$ & $0.446 \pm 0.062$ & $0.610 \pm 0.008$ \\
L25 & $0.130 \pm 0.305$ & $0.553 \pm 0.053$ & $0.479 \pm 0.039$ & $0.609 \pm 0.010$ \\
L26 & $0.251 \pm 0.254$ & $0.535 \pm 0.046$ & $0.411 \pm 0.098$ & $0.611 \pm 0.008$ \\
L27 & $0.099 \pm 0.278$ & $0.552 \pm 0.027$ & $0.366 \pm 0.197$ & $0.604 \pm 0.011$ \\
L28 & $0.099 \pm 0.307$ & $0.544 \pm 0.034$ & $0.358 \pm 0.218$ & $0.596 \pm 0.015$ \\
L29 & $0.088 \pm 0.275$ & $0.530 \pm 0.034$ & $0.356 \pm 0.271$ & $0.585 \pm 0.010$ \\
L30 & $0.220 \pm 0.234$ & $0.511 \pm 0.029$ & $0.424 \pm 0.104$ & $0.574 \pm 0.008$ \\
L31 & $0.069 \pm 0.284$ & $0.488 \pm 0.054$ & $0.332 \pm 0.262$ & $0.557 \pm 0.016$ \\
L32 & $0.584 \pm 0.011$ & $0.535 \pm 0.015$ & $0.410 \pm 0.185$ & $0.550 \pm 0.007$ \\
\bottomrule
\end{tabular}%
}
\end{table}

\paragraph{Equivalence testing.}
To formally test whether the cross-architecture peak $R^2$ values are statistically equivalent (rather than merely failing to differ), we apply the Two One-Sided Tests (TOST) procedure~\cite{lakens2017equivalence} with equivalence bound $\varepsilon = 0.05$.
The observed difference is $\Delta = 0.003$ (OpenVLA L9: $0.637 \pm 0.008$; TraceVLA L11: $0.634 \pm 0.009$).
Both one-sided tests reject at $p < 0.0001$; the 90\% confidence interval for the difference is $[-0.002, 0.009]$, falling entirely within $[-0.05, 0.05]$.
A conventional two-sample $t$-test yields $p = 0.318$ (not significant), illustrating the distinction between absence of evidence for a difference and evidence of equivalence: the TOST procedure provides the latter.

\subsection{Global Multiple Comparison Correction}
\label{sec:bonferroni_correction}

The paper reports 63 statistical tests spanning permutation tests, $t$-tests, bootstrap comparisons, TOST equivalence, and Spearman correlations.
We apply a Bonferroni correction ($\alpha_{\text{adj}} = 0.05 / 63 = 7.94 \times 10^{-4}$) and the less conservative Holm--Bonferroni step-down procedure; both yield identical results.
Of the 45 originally significant tests, 40 survive correction (88.9\%).

All eight core tests supporting the paper's main claims remain significant:
RSA permutation at L12 ($p < 0.0001$), trained vs.\ reference image geometry at L12 ($p = 3.9 \times 10^{-5}$), counterfactual wrong and empty instruction ($p < 10^{-100}$), TOST cross-architecture equivalence ($p < 0.0001$), dose-response top-24 ablation ($p < 10^{-100}$), trained vs.\ reference probe $R^2$ ($p < 0.0001$), and bottom-24 specificity ($p < 10^{-90}$).

Five auxiliary tests lose significance under correction:
(i)~RSA trained vs.\ reference preaction at L12 ($p = 0.005 \to p_{\text{adj}} = 0.284$),
(ii)~DINOv2 vs.\ SigLIP RSA ($p = 0.032 \to p_{\text{adj}} = 1.0$),
(iii)~fused backbone vs.\ SigLIP RSA ($p = 0.011 \to p_{\text{adj}} = 0.673$),
(iv)~CV exact paired permutation ($p = 0.016 \to p_{\text{adj}} = 0.984$),
and (v)~$R^2$ projector vs.\ VLA Spearman correlation ($p = 0.003 \to p_{\text{adj}} = 0.183$).
These concern auxiliary comparisons (encoder ranking, coefficient-of-variation ordering) rather than the core redistribution, probing, or counterfactual findings; no main-text conclusion requires revision.

\subsection{Mechanistic Interpretation}
Matched-size ablation ($d$: instruction $0.65$, random $0.60$, image $0.58$; Table~\ref{tab:specificity_control}) and dose-response (steep through 12 heads, plateau beyond 20) indicate distributed attention redistribution.
Counterfactual perturbation produces large effects across all conditions ($d = 1.0$--$1.2$ on Goal; \S\ref{sec:openvla_routing}), with shuffled tokens exceeding paraphrased ($d = 1.096$ vs.\ $0.996$) despite identical token sets, indicating positional sensitivity.
That ablation raises $R^2$ yet alters trajectories ($d = 0.65$) without increasing ground-truth error ($\Delta\text{MSE} \approx 0$; \S\ref{sec:action_trajectory}) reflects an information-computation dissociation~\cite{hewitt2019control}: instruction-dominant heads control direction rather than precision. Random ablation yields larger $R^2$ gains ($+0.071$ vs.\ $+0.030$; $z = -3.11$, $0$\textsuperscript{th} percentile), consistent with distinct $R^2$ signatures.

\subsection{Functional Stages and Predictions}
Per-layer profiles suggest three stages: redistribution (L0--L1; probe $99.6\%$, 24/32 instruction-dominant heads), computation (L8--L13; $R^2$ peak $0.655$--$0.678$), and decay (L14--L31, ${\sim}56\%$ of parameters; declining $R^2$ and RSA suggest pruning potential).
Task-specific fine-tuning eliminates L31 re-engagement (18 $\to$ 0 heads; \S\ref{sec:openvla_routing}), consistent with exit-layer attention serving multi-task flexibility ($N{=}2$).
TraceVLA~\cite{zheng2024tracevla} (Phi-3-V) shows consistent entry-exit attention concentration from an opposite starting point (\S\ref{sec:openvla_routing}).
The $2 {\times} 2$ trained-vs-reference comparison reveals a consistent reorganization pattern: reference OpenVLA collapses from L0 ($R^2 = 0.629$) to near-zero at mid-layers (L20: $0.086$, std $= 0.277$), while reference Phi-3-V peaks at L6 ($0.578$) with gradual decline (std $= 0.068$), producing qualitatively opposite reference landscapes.
Yet training drives both to nearly identical mid-layer peaks ($0.637$ vs.\ $0.634$, $\Delta = 0.003$) with stabilized deep-layer variance (OpenVLA std $0.277 \to 0.040$; Phi-3-V $0.068 \to 0.011$).
This consistent outcome from divergent starting points suggests that mid-layer action specialization is a consistent convergence pattern of VLA training, not an architectural artifact.
Exit-layer $R^2$ differs ($0.488$ vs.\ $0.557$), possibly reflecting architectural differences in late-layer processing.
Remaining predictions: cross-attention VLAs ($\pi_0$~\cite{black2024_pi0}, Octo~\cite{ghosh2024_octo}) should lack L1 instruction-dominant heads; larger VLAs should yield flatter per-head contributions.
For deployments, L1 instruction probes can flag instruction-discriminative tasks and L1 attention entropy can monitor instruction under-attendance at runtime.

\subsection{Extended Discussion}
\label{sec:extended_discussion_appendix}

\paragraph{Observational three-tier gradient.}
The cross-architecture matrix suggests a three-tier gradient of representational adaptation.
Tier~1 (VLM bases, no action training): inherited L1 attention (TraceVLA 32/32, NORA 13--14/16), negligible exit collapse ($\Delta R^2 \leq 0.048$), and minimal counterfactual sensitivity (cosine distance $\leq 0.185$).
Tier~2 (general-purpose VLA, OXE-pretrained): strong L1 routing (24--26/32 heads), moderate counterfactual sensitivity ($d \approx 1.0$), and no systematic exit collapse across suites.
Tier~3 (task-specific VLAs, fine-tuned on LIBERO): exit collapse at every VLM-to-VLA cell and medium-to-large counterfactual sensitivity ($d = 0.45$--$1.52$).
This gradient is suggestive but limited: Tier~2 contains a single architecture (OpenVLA), and the Tier~2/Tier~3 boundary confounds training stage with architecture.
Testing additional general-purpose VLAs (e.g., Octo~\cite{ghosh2024_octo}, $\pi_0$~\cite{black2024_pi0}) at Tier~2 is necessary before the three-tier structure can be treated as a robust empirical pattern.

\paragraph{Exit collapse mechanism.}
One possible explanation for the VLM-to-VLA / VLA-to-VLA asymmetry is that VLM-to-VLA references lack action-relevant exit-layer structure, forcing mid-layer construction without exit preservation (analogous to fine-tuning-induced feature distortion~\cite{kumar2022finetuning}), while VLA-to-VLA references already maintain this structure.
Validating this explanation causally requires controlled training-stage experiments, but the consistency across architectures supports its directional value.

\paragraph{Behavioral relevance.}
Correlating trained $\Delta R^2$ with published LIBERO success rates across all nine available architecture-suite cells yields a positive but non-significant trend (Spearman $\rho = 0.583$, $p = 0.10$, $n = 9$), marginally above the conventional significance threshold.
Within-suite comparisons, which control for task difficulty, are uniformly positive: on Goal ($n = 3$ architectures), $\Delta R^2$ and success rate are perfectly rank-correlated ($\rho = 1.0$), and every remaining suite ($n = 2$ each) shows the same direction.
This directional evidence is consistent with stronger mid-network action encoding (higher peak $R^2$) followed by more aggressive exit-layer compression being associated with better task performance.
Several caveats temper this interpretation: architecture capacity differences (model size, backbone quality) are not controlled, suite difficulty co-varies with both $\Delta R^2$ and success rate, multiple suites from the same architecture are not independent observations, and the sample size limits statistical power.
Post-hoc power analysis (Fisher $Z$-transform, two-tailed $\alpha = 0.05$) yields achieved power $= 0.37$ at $n = 9$; reaching $80\%$ power would require $n \geq 21$ architecture-suite cells.
The current analysis is therefore underpowered, and the non-significant $p$-value should not be interpreted as evidence of absence; it reflects insufficient sample size rather than a weak effect.
This correlational evidence does not establish a causal link between exit collapse and task success.
Closed-loop behavioral validation, including intervention experiments that directly manipulate exit-layer representations, remains a priority for future work.

\paragraph{Closed-loop evaluation attempt.}
We attempted closed-loop evaluation on LIBERO-Goal using the OpenVLA official checkpoint but achieved SR$=$0/3 across tested configurations, consistent with known reproducibility challenges in VLA deployment (sensitivity to hardware setup, action tokenization, and environment configuration).
This negative result motivates our representational analysis approach, which does not require successful task completion.

\paragraph{Metric constraints.}
Counterfactual sensitivity is measured via Cohen's $d$ on action MSE for models that produce actions and via cosine distance on hidden representations for reference VLMs (\S\ref{sec:cf_sensitivity}).
These metrics differ in scale and distributional properties; ratios between them are not meaningful.
The OpenVLA Long-Horizon inversion (trained $d = 0.56$, reference $d = 1.05$) illustrates how same-metric comparisons can reverse the expected direction, cautioning against overinterpretation of any single cell.
A standardized metric framework for cross-condition counterfactual comparison remains an open problem.

\paragraph{Extended limitations.}
Each architecture is analyzed at a single trained checkpoint, precluding characterization of how exit collapse or counterfactual sensitivity develop during training.
All analyses are open-loop, consistent with prior representational studies~\cite{grant2026,haon2025}; closed-loop validation would provide complementary evidence about the behavioral relevance of the identified routing structures.
The three-tier gradient has $n = 1$ at Tier~2, limiting its inferential power; populating this tier with additional general-purpose VLAs is a priority for future work.

\paragraph{Coarse training dynamics.}
Although intermediate checkpoints are unavailable (\S\ref{sec:limitations}), the cross-condition design provides coarse-grained evidence about when exit collapse emerges.
Figure~\ref{fig:training_dynamics} plots $\Delta R^2$ across four training-pathway categories.
NLP fine-tuning ($\Delta R^2 = 0.000$) and pretrained references ($\Delta R^2 \leq 0.097$) cluster below the $0.10$ threshold; VLM-to-VLA fine-tuning produces large collapse ($\Delta R^2 = 0.156$--$0.455$); subsequent VLA-to-VLA fine-tuning does not intensify it ($\Delta R^2 = 0.033$--$0.086$).
This constrains the onset: exit collapse emerges during the first action-grounded training stage and does not accumulate with continued task-specific fine-tuning.

\begin{figure}[h]
\centering
\includegraphics[width=0.55\columnwidth]{figures/fig_training_dynamics.pdf}
\caption{\textbf{$\Delta R^2$ across training pathways.} Each point is one architecture-suite cell; markers distinguish backbone families. Exit collapse (above dashed line) is confined to VLM-to-VLA transitions and does not intensify during VLA-to-VLA fine-tuning.}
\label{fig:training_dynamics}
\end{figure}

\subsection{Extended Practical Implications}
\label{sec:extended_practical_implications}

The main text identifies three use cases for the entry-exit decomposition; here we elaborate on each.

\paragraph{Checkpoint selection.}
$\Delta R^2$ provides a probe-based screening signal ($> 0.10$ flags reorganization) requiring only a forward pass and a linear probe.
Its wall-clock cost scales with model size rather than environment complexity, unlike rollout evaluation which requires physics simulation and multi-step execution.
Practitioners can rank checkpoints by $\Delta R^2$ and reserve rollout budgets for the top candidates, substantially reducing the number of expensive closed-loop evaluations needed during model selection.

\paragraph{Transfer learning strategy.}
The flat VLA-to-VLA $\Delta R^2$ profile (\S\ref{sec:r2_collapse}) shows that VLA-to-VLA transitions preserve exit structure, justifying staged pipelines (VLM$\to$VLA$\to$task) that reuse the first-stage action head without re-learning exit representations.
Inherited L1 routing can further serve as a frozen prior, reducing tunable parameters during domain transfer.
This suggests that practitioners adapting a general-purpose VLA to a new task suite may benefit from freezing early layers while focusing capacity on mid-to-late layer adaptation.

\paragraph{Mid-layer probe value.}
Where exit collapse is pronounced, mid-layer representations ($R^2$ peak at L9--L11 across architectures) outperform exit-layer features for downstream tasks such as reward modeling, distillation targets, or auxiliary prediction heads.
For knowledge distillation, matching mid-layer features of the teacher rather than exit-layer features may yield better student performance in VLA settings.
For reward modeling, training reward predictors on mid-layer features rather than final-layer outputs may provide more informative training signals.

\subsection{Methodological Context}
Representational analysis is an established methodology in the interpretability tradition; behavioral validation provides complementary evidence but is not a prerequisite for mechanistic characterization.
Table~\ref{tab:methodological_precedent} situates our approach within this tradition.

\begin{table}[h]
\centering
\caption{\textbf{Methodological precedent.} Representative interpretability studies and their use of behavioral evaluation. Partial behavioral eval includes correlational or indirect evidence.}
\label{tab:methodological_precedent}
\small
\begin{tabular}{@{}llll@{}}
\toprule
\textbf{Paper} & \textbf{Domain} & \textbf{Method} & \textbf{Behav.\ eval?} \\
\midrule
Belinkov~\cite{belinkov2022} & NLP & Probing survey & No \\
Olsson et al.~\cite{olsson2022} & LM & Induction heads & No \\
Wang et al.~\cite{wang2023_ioi} & LM & IOI circuit & No \\
Meng et al.~\cite{meng2022_rome} & LM & Causal tracing (localization) & No \\
Grant et al.~\cite{grant2026} & VLA & Visual dominance & Partial \\
Haon et al.~\cite{haon2025} & VLA & Sparse directions & No \\
\midrule
Ours & VLA & $\Delta R^2$ + attn.\ profiling & Partial ($\rho{=}0.583$) \\
\bottomrule
\end{tabular}
\end{table}

\subsection{NLP Finetuning Controls}
\label{sec:nlp_control_appendix}

Table~\ref{tab:nlp_controls} compares the three NLP finetuning controls used to test whether exit collapse is specific to action-oriented training.
All three chat-finetuned models produce flat action-probe profiles with $\Delta R^2 = 0.000$, confirming that NLP instruction tuning does not induce the peak-to-exit drop observed in VLM-to-VLA transitions in any of the three backbone families.

\begin{table}[h]
\centering
\caption{\textbf{NLP finetuning control comparison.} All three models are instruction-tuned for conversational use without action heads or robot demonstration data. None exhibits exit collapse.}
\label{tab:nlp_controls}
\small
\begin{tabular}{@{}llccc@{}}
\toprule
\textbf{Model} & \textbf{VLA backbone} & \textbf{$\Delta R^2$} & \textbf{Mean $R^2$} & \textbf{Layers} \\
\midrule
Llama-2-7b-chat & OpenVLA (Llama-2-7B) & 0.000 & $0.155 \pm 0.001$ & 32 \\
Phi-3-mini-4k-instruct & TraceVLA (Phi-3-mini) & 0.000 & $0.155 \pm 0.000$ & 32 \\
Qwen2.5-3B-Instruct$^{*}$ & NORA (Qwen2.5-VL-3B) & 0.000 & $0.155 \pm 0.000$ & 36 \\
\bottomrule
\multicolumn{5}{l}{\footnotesize $^{*}$Qwen2.5-3B-Instruct shares the Qwen2.5 language backbone with NORA's}\\
\multicolumn{5}{l}{\footnotesize Qwen2.5-VL-3B but lacks the vision encoder and cross-attention modules.}\\
\end{tabular}
\end{table}

Figure~\ref{fig:r2_collapse} (main text) overlays the NLP control profiles (dotted gray) against the corresponding VLA trained and reference curves.

\subsection{Claim-Evidence Summary}
\label{sec:claim_evidence_appendix}

\begin{figure}[h]
\centering
\includegraphics[width=\columnwidth]{figures/fig_claim_evidence.pdf}
\caption{\textbf{Claim--evidence--scope summary.} Each main finding is paired with its supporting analysis and its main limitation. The figure distinguishes directly observed representational patterns from stronger causal interpretations that require interventional or training-dynamics evidence.}
\end{figure}

\subsection{Cross-Architecture Design Rationale}
The two VLM-to-VLA transitions (TraceVLA and NORA) were selected to maximize structural diversity within the constraint of publicly available LIBERO-evaluated checkpoints.
Table~\ref{tab:arch_diversity} summarizes the design-axis differences.

\begin{table}[h]
\centering
\caption{\textbf{Structural diversity between the two VLM-to-VLA transitions.} The two architectures differ on nearly every major design axis, strengthening the generality of shared findings.}
\label{tab:arch_diversity}
\small
\begin{tabular}{@{}lll@{}}
\toprule
\textbf{Design Axis} & \textbf{TraceVLA} & \textbf{NORA} \\
\midrule
Language backbone & Phi-3-mini (3.8B) & Qwen2.5-VL (3B) \\
Backbone family & Microsoft & Alibaba \\
Attention mechanism & Multi-Head Attention & Grouped-Query Attention \\
Number of layers & 32 & 36 \\
Vision encoder & SigLIP (shared w/ base) & Native ViT (dynamic resolution) \\
Vision integration & Projector (late fusion) & Token concatenation (native) \\
Pretraining data & English-centric web & Multilingual web + code \\
Action head design & Linear projection & MLP \\
Training data & OXE + LIBERO & OXE + LIBERO + custom \\
\bottomrule
\end{tabular}
\end{table}

Despite these differences, both transitions exhibit the same exit-collapse pattern ($\Delta R^2 = 0.156$--$0.455$), while their respective reference VLMs show only small fluctuations ($\Delta R^2 \leq 0.097$).
The probability of observing consistent collapse in two structurally divergent architectures under a null hypothesis of random association is low (permutation test $p = 0.008$; \S\ref{sec:delta_r2_threshold}).
Structural diversity thus provides evidential weight beyond what the raw $N = 2$ count suggests.

\subsection{Convergent Evidence Summary}
Three methodologically distinct approaches triangulate instruction routing in VLAs.
Table~\ref{tab:convergent_evidence} summarizes what each method measures, its key finding, evidence type, and acknowledged limitation.

\begin{table}[h]
\centering
\caption{\textbf{Convergent evidence across three analytical methods.} Each method addresses a different aspect of instruction routing; no single method is sufficient alone, but their convergence strengthens the overall characterization.}
\label{tab:convergent_evidence}
\scriptsize
\setlength{\tabcolsep}{3pt}
\begin{tabular}{@{}p{2.0cm}p{2.8cm}p{3.2cm}p{1.8cm}p{2.8cm}@{}}
\toprule
\textbf{Method} & \textbf{What it measures} & \textbf{Key finding} & \textbf{Evidence type} & \textbf{Limitation} \\
\midrule
Layer-wise probing (\S\ref{sec:r2_collapse}) & Where action-relevant information is decodable & VLM-to-VLA: peak-to-exit collapse ($\Delta R^2 = 0.16$--$0.46$); VLA-to-VLA: flat & Correlational (decodability) & Does not establish causal necessity \\
\addlinespace
Counterfactual perturbation (\S\ref{sec:cf_sensitivity}) & Whether instruction content affects model outputs & All trained VLAs: graded divergence ($d = 0.45$--$1.52$) ordered by perturbation severity & Behavioral (input--output) & Does not localize which layers are responsible \\
\addlinespace
Attention head ablation (Appendix~\ref{sec:action_trajectory}) & Whether identified structures carry functional signals & Dose-response: monotonic divergence ($d = 0.48$--$0.64$) upon progressive head removal & Interventional (causal contribution) & Single architecture (OpenVLA); contribution $\neq$ necessity \\
\bottomrule
\end{tabular}
\end{table}

The three methods form a complementary chain: probing identifies \emph{where} instruction information concentrates across layers, perturbation confirms \emph{whether} instruction content functionally affects behavior, and ablation demonstrates \emph{that} the identified attention structures carry relevant signals.
No method alone supports causal claims about instruction routing mechanisms; their convergence provides stronger evidence than any single approach.

\section{Use of Generative AI}
\label{sec:use_of_ai}

Generative AI tools (Claude, GPT-4) assisted with code implementation and manuscript editing; all scientific content was produced by the authors.

